\documentclass[11pt]{article}
\usepackage[margin=0.85in]{geometry}
\usepackage[T1]{fontenc}
\usepackage{lmodern}
\usepackage{amsmath,amssymb,amsthm,mathtools,booktabs,longtable,array}
\usepackage{listings,xcolor,xurl,hyperref}
\setlength{\emergencystretch}{2em}
\hypersetup{colorlinks=true,linkcolor=blue,urlcolor=blue,citecolor=blue,
 pdftitle={A truncated distance bound for cyclic lifted-product quantum codes},pdfauthor={}}
\newtheorem{theorem}{Theorem}[section]
\newtheorem{lemma}[theorem]{Lemma}
\newtheorem{proposition}[theorem]{Proposition}
\theoremstyle{definition}\newtheorem{definition}[theorem]{Definition}
\newcommand{\F}{\mathbb F_2}\newcommand{\Z}{\mathbb Z}
\newcommand{\wt}{\operatorname{wt}}\newcommand{\row}{\operatorname{row}}
\newcommand{\HX}{H_{\mathsf X}}\newcommand{\HZ}{H_{\mathsf Z}}
\newcommand{\supp}{\operatorname{supp}}
\newcommand{\lean}[2]{\csname proof@#1.#2\endcsname}
\newcolumntype{P}[1]{>{\raggedright\arraybackslash}p{#1}}
\expandafter\def\csname proof@CoverageCore.Covered\endcsname{\href{https://github.com/alectionary/vibe-math/blob/3ae21c0719cedbd928f4f7882f9b9e266567605e/formalization/distance-preservation/truncated12/CoverageCore.lean\#L7-L8}{\nolinkurl{KernelCatalog.Covered}}}
\expandafter\def\csname proof@CoverageCore.covered_single\endcsname{\href{https://github.com/alectionary/vibe-math/blob/3ae21c0719cedbd928f4f7882f9b9e266567605e/formalization/distance-preservation/truncated12/CoverageCore.lean\#L10-L15}{\nolinkurl{KernelCatalog.covered_single}}}
\expandafter\def\csname proof@CoverageCore.covered_append\endcsname{\href{https://github.com/alectionary/vibe-math/blob/3ae21c0719cedbd928f4f7882f9b9e266567605e/formalization/distance-preservation/truncated12/CoverageCore.lean\#L17-L26}{\nolinkurl{KernelCatalog.covered_append}}}
\expandafter\def\csname proof@Coverage10.complete_coverage10\endcsname{\href{https://github.com/alectionary/vibe-math/blob/3ae21c0719cedbd928f4f7882f9b9e266567605e/formalization/distance-preservation/truncated12/Coverage10.lean\#L248968-L249217}{\nolinkurl{KernelCatalog.complete_coverage10}}}
\expandafter\def\csname proof@GeneralLP.HX\endcsname{\href{https://github.com/alectionary/vibe-math/blob/3ae21c0719cedbd928f4f7882f9b9e266567605e/formalization/distance-preservation/GeneralLP.lean\#L11-L13}{\nolinkurl{GeneralLP.HX}}}
\expandafter\def\csname proof@GeneralLP.HZ\endcsname{\href{https://github.com/alectionary/vibe-math/blob/3ae21c0719cedbd928f4f7882f9b9e266567605e/formalization/distance-preservation/GeneralLP.lean\#L15-L17}{\nolinkurl{GeneralLP.HZ}}}
\expandafter\def\csname proof@GeneralLP.InSpanX\endcsname{\href{https://github.com/alectionary/vibe-math/blob/3ae21c0719cedbd928f4f7882f9b9e266567605e/formalization/distance-preservation/GeneralLP.lean\#L65-L67}{\nolinkurl{GeneralLP.InSpanX}}}
\expandafter\def\csname proof@ActualIncidence.classical_incidence\endcsname{\href{https://github.com/alectionary/vibe-math/blob/3ae21c0719cedbd928f4f7882f9b9e266567605e/formalization/distance-preservation/truncated12/ActualIncidence.lean\#L19-L22}{\nolinkurl{TruncatedDistance.classical_incidence}}}
\expandafter\def\csname proof@ActualIncidence.quantum_incidence\endcsname{\href{https://github.com/alectionary/vibe-math/blob/3ae21c0719cedbd928f4f7882f9b9e266567605e/formalization/distance-preservation/truncated12/ActualIncidence.lean\#L24-L33}{\nolinkurl{TruncatedDistance.quantum_incidence}}}
\expandafter\def\csname proof@ActualIncidence.quantum_no_four\endcsname{\href{https://github.com/alectionary/vibe-math/blob/3ae21c0719cedbd928f4f7882f9b9e266567605e/formalization/distance-preservation/truncated12/ActualIncidence.lean\#L64-L77}{\nolinkurl{TruncatedDistance.quantum_no_four}}}
\expandafter\def\csname proof@ActualIncidence.quantum_no_six\endcsname{\href{https://github.com/alectionary/vibe-math/blob/3ae21c0719cedbd928f4f7882f9b9e266567605e/formalization/distance-preservation/truncated12/ActualIncidence.lean\#L79-L111}{\nolinkurl{TruncatedDistance.quantum_no_six}}}
\expandafter\def\csname proof@ActualIncidence.source_kernel_zero_or_two\endcsname{\href{https://github.com/alectionary/vibe-math/blob/3ae21c0719cedbd928f4f7882f9b9e266567605e/formalization/distance-preservation/truncated12/ActualIncidence.lean\#L127-L144}{\nolinkurl{TruncatedDistance.source_kernel_zero_or_two}}}
\expandafter\def\csname proof@BlockParity.common_block_parity\endcsname{\href{https://github.com/alectionary/vibe-math/blob/3ae21c0719cedbd928f4f7882f9b9e266567605e/formalization/distance-preservation/truncated12/BlockParity.lean\#L27-L40}{\nolinkurl{TruncatedDistance.common_block_parity}}}
\expandafter\def\csname proof@BlockCounts.common_count_parity\endcsname{\href{https://github.com/alectionary/vibe-math/blob/3ae21c0719cedbd928f4f7882f9b9e266567605e/formalization/distance-preservation/truncated12/BlockCounts.lean\#L36-L48}{\nolinkurl{TruncatedDistance.common_count_parity}}}
\expandafter\def\csname proof@TorsionClassical.torsion_card_le\endcsname{\href{https://github.com/alectionary/vibe-math/blob/3ae21c0719cedbd928f4f7882f9b9e266567605e/formalization/distance-preservation/truncated12/TorsionClassical.lean\#L10-L13}{\nolinkurl{TruncatedDistance.torsion_card_le}}}
\expandafter\def\csname proof@TorsionClassical.torsionWord\endcsname{\href{https://github.com/alectionary/vibe-math/blob/3ae21c0719cedbd928f4f7882f9b9e266567605e/formalization/distance-preservation/truncated12/TorsionClassical.lean\#L17-L19}{\nolinkurl{TruncatedDistance.torsionWord}}}
\expandafter\def\csname proof@TorsionClassical.torsionWord_source_kernel\endcsname{\href{https://github.com/alectionary/vibe-math/blob/3ae21c0719cedbd928f4f7882f9b9e266567605e/formalization/distance-preservation/truncated12/TorsionClassical.lean\#L54-L64}{\nolinkurl{TruncatedDistance.torsionWord_source_kernel}}}
\expandafter\def\csname proof@TorsionClassical.source_classical_word_le10\endcsname{\href{https://github.com/alectionary/vibe-math/blob/3ae21c0719cedbd928f4f7882f9b9e266567605e/formalization/distance-preservation/truncated12/TorsionClassical.lean\#L66-L74}{\nolinkurl{TruncatedDistance.source_classical_word_le10}}}
\expandafter\def\csname proof@IntegerSoundness.integer_certificate_sound\endcsname{\href{https://github.com/alectionary/vibe-math/blob/3ae21c0719cedbd928f4f7882f9b9e266567605e/formalization/distance-preservation/truncated12/IntegerSoundness.lean\#L11-L22}{\nolinkurl{TruncatedDistance.integer_certificate_sound}}}
\expandafter\def\csname proof@IntegerSoundness.row_small_torsion_impossible\endcsname{\href{https://github.com/alectionary/vibe-math/blob/3ae21c0719cedbd928f4f7882f9b9e266567605e/formalization/distance-preservation/truncated12/IntegerSoundness.lean\#L61-L65}{\nolinkurl{TruncatedDistance.row_small_torsion_impossible}}}
\expandafter\def\csname proof@IntegerSoundness.column_two_torsion_impossible\endcsname{\href{https://github.com/alectionary/vibe-math/blob/3ae21c0719cedbd928f4f7882f9b9e266567605e/formalization/distance-preservation/truncated12/IntegerSoundness.lean\#L67-L71}{\nolinkurl{TruncatedDistance.column_two_torsion_impossible}}}
\expandafter\def\csname proof@SupportSingletons.actual_singleton_label_alternative\endcsname{\href{https://github.com/alectionary/vibe-math/blob/3ae21c0719cedbd928f4f7882f9b9e266567605e/formalization/distance-preservation/truncated12/SupportSingletons.lean\#L71-L90}{\nolinkurl{TruncatedDistance.actual_singleton_label_alternative}}}
\expandafter\def\csname proof@KernelCatalog.checked_fast_equations\endcsname{\href{https://github.com/alectionary/vibe-math/blob/3ae21c0719cedbd928f4f7882f9b9e266567605e/formalization/distance-preservation/truncated12/KernelCatalog.lean\#L41-L57}{\nolinkurl{KernelCatalog.checked_fast_equations}}}
\expandafter\def\csname proof@PureDistance.pure_left_zero_below_classical\endcsname{\href{https://github.com/alectionary/vibe-math/blob/3ae21c0719cedbd928f4f7882f9b9e266567605e/formalization/distance-preservation/truncated12/PureDistance.lean\#L18-L30}{\nolinkurl{TruncatedDistance.pure_left_zero_below_classical}}}
\expandafter\def\csname proof@PureDistance.pure_right_zero_below_transpose\endcsname{\href{https://github.com/alectionary/vibe-math/blob/3ae21c0719cedbd928f4f7882f9b9e266567605e/formalization/distance-preservation/truncated12/PureDistance.lean\#L32-L44}{\nolinkurl{TruncatedDistance.pure_right_zero_below_transpose}}}
\expandafter\def\csname proof@TargetObstructions.source_target_obstruction\endcsname{\href{https://github.com/alectionary/vibe-math/blob/3ae21c0719cedbd928f4f7882f9b9e266567605e/formalization/distance-preservation/truncated12/TargetObstructions.lean\#L1695-L1767}{\nolinkurl{TruncatedDistance.source_target_obstruction}}}
\expandafter\def\csname proof@PatternRealizations.PatternRealization\endcsname{\href{https://github.com/alectionary/vibe-math/blob/3ae21c0719cedbd928f4f7882f9b9e266567605e/formalization/distance-preservation/truncated12/PatternRealizations.lean\#L10-L21}{\nolinkurl{TruncatedDistance.PatternRealization}}}
\expandafter\def\csname proof@PatternRealizations.certified_pattern_has_no_source_realization\endcsname{\href{https://github.com/alectionary/vibe-math/blob/3ae21c0719cedbd928f4f7882f9b9e266567605e/formalization/distance-preservation/truncated12/PatternRealizations.lean\#L54-L72}{\nolinkurl{TruncatedDistance.certified_pattern_has_no_source_realization}}}
\expandafter\def\csname proof@GeneralSymmetry.fullSpanX\endcsname{\href{https://github.com/alectionary/vibe-math/blob/3ae21c0719cedbd928f4f7882f9b9e266567605e/formalization/distance-preservation/truncated12/GeneralSymmetry.lean\#L22-L23}{\nolinkurl{TruncatedDistance.fullSpanX}}}
\expandafter\def\csname proof@GeneralSymmetry.fullSpanZ\endcsname{\href{https://github.com/alectionary/vibe-math/blob/3ae21c0719cedbd928f4f7882f9b9e266567605e/formalization/distance-preservation/truncated12/GeneralSymmetry.lean\#L24-L25}{\nolinkurl{TruncatedDistance.fullSpanZ}}}
\expandafter\def\csname proof@GeneralSymmetry.full_span_swap\endcsname{\href{https://github.com/alectionary/vibe-math/blob/3ae21c0719cedbd928f4f7882f9b9e266567605e/formalization/distance-preservation/truncated12/GeneralSymmetry.lean\#L27-L43}{\nolinkurl{TruncatedDistance.full_span_swap}}}
\expandafter\def\csname proof@GeneralSymmetry.kernel_swap\endcsname{\href{https://github.com/alectionary/vibe-math/blob/3ae21c0719cedbd928f4f7882f9b9e266567605e/formalization/distance-preservation/truncated12/GeneralSymmetry.lean\#L45-L60}{\nolinkurl{TruncatedDistance.kernel_swap}}}
\expandafter\def\csname proof@GeneralSymmetry.swap_weight\endcsname{\href{https://github.com/alectionary/vibe-math/blob/3ae21c0719cedbd928f4f7882f9b9e266567605e/formalization/distance-preservation/truncated12/GeneralSymmetry.lean\#L62-L65}{\nolinkurl{TruncatedDistance.swap_weight}}}
\expandafter\def\csname proof@RowCancellation.single_exception_cancellation\endcsname{\href{https://github.com/alectionary/vibe-math/blob/3ae21c0719cedbd928f4f7882f9b9e266567605e/formalization/distance-preservation/truncated12/RowCancellation.lean\#L36-L85}{\nolinkurl{TruncatedDistance.single_exception_cancellation}}}
\expandafter\def\csname proof@ConstantRows.constant_singleton_rows\endcsname{\href{https://github.com/alectionary/vibe-math/blob/3ae21c0719cedbd928f4f7882f9b9e266567605e/formalization/distance-preservation/truncated12/ConstantRows.lean\#L25-L43}{\nolinkurl{TruncatedDistance.constant_singleton_rows}}}
\expandafter\def\csname proof@SourceSeven.seven_certificates_complete\endcsname{\href{https://github.com/alectionary/vibe-math/blob/3ae21c0719cedbd928f4f7882f9b9e266567605e/formalization/distance-preservation/truncated12/SourceSeven.lean\#L12-L158}{\nolinkurl{TruncatedDistance.seven_certificates_complete}}}
\expandafter\def\csname proof@SourceSeven.source_seven_singleton_profile_stabilizer\endcsname{\href{https://github.com/alectionary/vibe-math/blob/3ae21c0719cedbd928f4f7882f9b9e266567605e/formalization/distance-preservation/truncated12/SourceSeven.lean\#L258-L281}{\nolinkurl{TruncatedDistance.source_seven_singleton_profile_stabilizer}}}
\expandafter\def\csname proof@LocalMantel.source_left_local_mantel\endcsname{\href{https://github.com/alectionary/vibe-math/blob/3ae21c0719cedbd928f4f7882f9b9e266567605e/formalization/distance-preservation/truncated12/LocalMantel.lean\#L39-L57}{\nolinkurl{TruncatedDistance.source_left_local_mantel}}}
\expandafter\def\csname proof@LocalMantel.source_right_local_mantel\endcsname{\href{https://github.com/alectionary/vibe-math/blob/3ae21c0719cedbd928f4f7882f9b9e266567605e/formalization/distance-preservation/truncated12/LocalMantel.lean\#L59-L77}{\nolinkurl{TruncatedDistance.source_right_local_mantel}}}
\expandafter\def\csname proof@CountPatterns.count_cases_length\endcsname{\href{https://github.com/alectionary/vibe-math/blob/3ae21c0719cedbd928f4f7882f9b9e266567605e/formalization/distance-preservation/truncated12/CountPatterns.lean\#L69-L69}{\nolinkurl{TruncatedDistance.count_cases_length}}}
\expandafter\def\csname proof@CountPatterns.classified_of_mem_countCases\endcsname{\href{https://github.com/alectionary/vibe-math/blob/3ae21c0719cedbd928f4f7882f9b9e266567605e/formalization/distance-preservation/truncated12/CountPatterns.lean\#L77-L82}{\nolinkurl{TruncatedDistance.classified_of_mem_countCases}}}
\expandafter\def\csname proof@SourceProfiles.source_mixed_count_profiles\endcsname{\href{https://github.com/alectionary/vibe-math/blob/3ae21c0719cedbd928f4f7882f9b9e266567605e/formalization/distance-preservation/truncated12/SourceProfiles.lean\#L29-L79}{\nolinkurl{TruncatedDistance.source_mixed_count_profiles}}}
\expandafter\def\csname proof@PauliBridge.pauli_bound_of_x_bound\endcsname{\href{https://github.com/alectionary/vibe-math/blob/3ae21c0719cedbd928f4f7882f9b9e266567605e/formalization/distance-preservation/truncated12/PauliBridge.lean\#L31-L42}{\nolinkurl{TruncatedDistance.pauli_bound_of_x_bound}}}
\expandafter\def\csname proof@TenGeometry.source_ten_geometry\endcsname{\href{https://github.com/alectionary/vibe-math/blob/3ae21c0719cedbd928f4f7882f9b9e266567605e/formalization/distance-preservation/truncated12/TenGeometry.lean\#L26-L87}{\nolinkurl{TruncatedDistance.source_ten_geometry}}}
\expandafter\def\csname proof@TenIndex.tenRank_bound\endcsname{\href{https://github.com/alectionary/vibe-math/blob/3ae21c0719cedbd928f4f7882f9b9e266567605e/formalization/distance-preservation/truncated12/TenIndex.lean\#L10-L15}{\nolinkurl{TruncatedDistance.tenRank_bound}}}
\expandafter\def\csname proof@TenEncoding.ten_geometry_encodes\endcsname{\href{https://github.com/alectionary/vibe-math/blob/3ae21c0719cedbd928f4f7882f9b9e266567605e/formalization/distance-preservation/truncated12/TenEncoding.lean\#L8-L38}{\nolinkurl{TruncatedDistance.ten_geometry_encodes}}}
\expandafter\def\csname proof@TenEncoding.actual_canonical_ten_support_encodes\endcsname{\href{https://github.com/alectionary/vibe-math/blob/3ae21c0719cedbd928f4f7882f9b9e266567605e/formalization/distance-preservation/truncated12/TenEncoding.lean\#L40-L52}{\nolinkurl{TruncatedDistance.actual_canonical_ten_support_encodes}}}
\expandafter\def\csname proof@ExceptionalGraph.exceptional_group_is_k23\endcsname{\href{https://github.com/alectionary/vibe-math/blob/3ae21c0719cedbd928f4f7882f9b9e266567605e/formalization/distance-preservation/truncated12/ExceptionalGraph.lean\#L25-L51}{\nolinkurl{TruncatedDistance.exceptional_group_is_k23}}}
\expandafter\def\csname proof@K23Numbering.actual_k23_numbering\endcsname{\href{https://github.com/alectionary/vibe-math/blob/3ae21c0719cedbd928f4f7882f9b9e266567605e/formalization/distance-preservation/truncated12/K23Numbering.lean\#L29-L74}{\nolinkurl{TruncatedDistance.actual_k23_numbering}}}
\expandafter\def\csname proof@ElevenGeometry.actual_eleven_geometry\endcsname{\href{https://github.com/alectionary/vibe-math/blob/3ae21c0719cedbd928f4f7882f9b9e266567605e/formalization/distance-preservation/truncated12/ElevenGeometry.lean\#L74-L115}{\nolinkurl{TruncatedDistance.actual_eleven_geometry}}}
\expandafter\def\csname proof@ElevenEncoding.eleven_geometry_encodes\endcsname{\href{https://github.com/alectionary/vibe-math/blob/3ae21c0719cedbd928f4f7882f9b9e266567605e/formalization/distance-preservation/truncated12/ElevenEncoding.lean\#L7-L18}{\nolinkurl{TruncatedDistance.eleven_geometry_encodes}}}
\expandafter\def\csname proof@Coverage11.complete_coverage11\endcsname{\href{https://github.com/alectionary/vibe-math/blob/3ae21c0719cedbd928f4f7882f9b9e266567605e/formalization/distance-preservation/truncated12/Coverage11.lean\#L117574-L117691}{\nolinkurl{KernelCatalog.complete_coverage11}}}
\expandafter\def\csname proof@SourceEleven.source_canonical_eleven_stabilizer\endcsname{\href{https://github.com/alectionary/vibe-math/blob/3ae21c0719cedbd928f4f7882f9b9e266567605e/formalization/distance-preservation/truncated12/SourceEleven.lean\#L15-L51}{\nolinkurl{TruncatedDistance.source_canonical_eleven_stabilizer}}}
\expandafter\def\csname proof@ProfileTransport.DistanceCriterion\endcsname{\href{https://github.com/alectionary/vibe-math/blob/3ae21c0719cedbd928f4f7882f9b9e266567605e/formalization/distance-preservation/truncated12/ProfileTransport.lean\#L9-L17}{\nolinkurl{TruncatedDistance.DistanceCriterion}}}
\expandafter\def\csname proof@ProfileTransport.normPerm_criterion\endcsname{\href{https://github.com/alectionary/vibe-math/blob/3ae21c0719cedbd928f4f7882f9b9e266567605e/formalization/distance-preservation/truncated12/ProfileTransport.lean\#L30-L43}{\nolinkurl{TruncatedDistance.normPerm_criterion}}}
\expandafter\def\csname proof@ProfileTransport.normPerm_kernel\endcsname{\href{https://github.com/alectionary/vibe-math/blob/3ae21c0719cedbd928f4f7882f9b9e266567605e/formalization/distance-preservation/truncated12/ProfileTransport.lean\#L45-L54}{\nolinkurl{TruncatedDistance.normPerm_kernel}}}
\expandafter\def\csname proof@ProfileTransport.normPerm_span\endcsname{\href{https://github.com/alectionary/vibe-math/blob/3ae21c0719cedbd928f4f7882f9b9e266567605e/formalization/distance-preservation/truncated12/ProfileTransport.lean\#L56-L64}{\nolinkurl{TruncatedDistance.normPerm_span}}}
\expandafter\def\csname proof@ProfileTransport.normPerm_weight\endcsname{\href{https://github.com/alectionary/vibe-math/blob/3ae21c0719cedbd928f4f7882f9b9e266567605e/formalization/distance-preservation/truncated12/ProfileTransport.lean\#L66-L73}{\nolinkurl{TruncatedDistance.normPerm_weight}}}
\expandafter\def\csname proof@CountCanonical.ten_left_marker\endcsname{\href{https://github.com/alectionary/vibe-math/blob/3ae21c0719cedbd928f4f7882f9b9e266567605e/formalization/distance-preservation/truncated12/CountCanonical.lean\#L32-L40}{\nolinkurl{TruncatedDistance.ten_left_marker}}}
\expandafter\def\csname proof@CountCanonical.ten_right_marker\endcsname{\href{https://github.com/alectionary/vibe-math/blob/3ae21c0719cedbd928f4f7882f9b9e266567605e/formalization/distance-preservation/truncated12/CountCanonical.lean\#L42-L50}{\nolinkurl{TruncatedDistance.ten_right_marker}}}
\expandafter\def\csname proof@CountCanonical.eleven_left_marker\endcsname{\href{https://github.com/alectionary/vibe-math/blob/3ae21c0719cedbd928f4f7882f9b9e266567605e/formalization/distance-preservation/truncated12/CountCanonical.lean\#L52-L68}{\nolinkurl{TruncatedDistance.eleven_left_marker}}}
\expandafter\def\csname proof@TwelveConditional.source_kernel_below_twelve_of_coverage\endcsname{\href{https://github.com/alectionary/vibe-math/blob/3ae21c0719cedbd928f4f7882f9b9e266567605e/formalization/distance-preservation/truncated12/TwelveConditional.lean\#L13-L133}{\nolinkurl{TruncatedDistance.source_kernel_below_twelve_of_coverage}}}
\expandafter\def\csname proof@ThroughEleven.source_kernel_below_twelve_is_stabilizer\endcsname{\href{https://github.com/alectionary/vibe-math/blob/3ae21c0719cedbd928f4f7882f9b9e266567605e/formalization/distance-preservation/truncated12/ThroughEleven.lean\#L10-L17}{\nolinkurl{TruncatedDistance.source_kernel_below_twelve_is_stabilizer}}}
\expandafter\def\csname proof@ThroughEleven.source_pauli_bound_twelve\endcsname{\href{https://github.com/alectionary/vibe-math/blob/3ae21c0719cedbd928f4f7882f9b9e266567605e/formalization/distance-preservation/truncated12/ThroughEleven.lean\#L27-L39}{\nolinkurl{TruncatedDistance.source_pauli_bound_twelve}}}
\expandafter\def\csname proof@DistanceMinima.sourceClassicalHasWeight\endcsname{\href{https://github.com/alectionary/vibe-math/blob/3ae21c0719cedbd928f4f7882f9b9e266567605e/formalization/distance-preservation/truncated12/DistanceMinima.lean\#L9-L11}{\nolinkurl{TruncatedDistance.sourceClassicalHasWeight}}}
\expandafter\def\csname proof@DistanceMinima.sourceClassicalDistance\endcsname{\href{https://github.com/alectionary/vibe-math/blob/3ae21c0719cedbd928f4f7882f9b9e266567605e/formalization/distance-preservation/truncated12/DistanceMinima.lean\#L21-L22}{\nolinkurl{TruncatedDistance.sourceClassicalDistance}}}
\expandafter\def\csname proof@DistanceMinima.source_classical_minimum_realized\endcsname{\href{https://github.com/alectionary/vibe-math/blob/3ae21c0719cedbd928f4f7882f9b9e266567605e/formalization/distance-preservation/truncated12/DistanceMinima.lean\#L24-L26}{\nolinkurl{TruncatedDistance.source_classical_minimum_realized}}}
\expandafter\def\csname proof@DistanceMinima.source_classical_minimum_lower\endcsname{\href{https://github.com/alectionary/vibe-math/blob/3ae21c0719cedbd928f4f7882f9b9e266567605e/formalization/distance-preservation/truncated12/DistanceMinima.lean\#L28-L31}{\nolinkurl{TruncatedDistance.source_classical_minimum_lower}}}
\expandafter\def\csname proof@DistanceMinima.sourceTransposeHasWeight\endcsname{\href{https://github.com/alectionary/vibe-math/blob/3ae21c0719cedbd928f4f7882f9b9e266567605e/formalization/distance-preservation/truncated12/DistanceMinima.lean\#L33-L35}{\nolinkurl{TruncatedDistance.sourceTransposeHasWeight}}}
\expandafter\def\csname proof@DistanceMinima.sourceTransposeDistance\endcsname{\href{https://github.com/alectionary/vibe-math/blob/3ae21c0719cedbd928f4f7882f9b9e266567605e/formalization/distance-preservation/truncated12/DistanceMinima.lean\#L37-L38}{\nolinkurl{TruncatedDistance.sourceTransposeDistance}}}
\expandafter\def\csname proof@DistanceMinima.source_transpose_zero_kernel_infinity\endcsname{\href{https://github.com/alectionary/vibe-math/blob/3ae21c0719cedbd928f4f7882f9b9e266567605e/formalization/distance-preservation/truncated12/DistanceMinima.lean\#L40-L51}{\nolinkurl{TruncatedDistance.source_transpose_zero_kernel_infinity}}}
\expandafter\def\csname proof@DistanceMinima.sourceQuantumHasWeight\endcsname{\href{https://github.com/alectionary/vibe-math/blob/3ae21c0719cedbd928f4f7882f9b9e266567605e/formalization/distance-preservation/truncated12/DistanceMinima.lean\#L60-L65}{\nolinkurl{TruncatedDistance.sourceQuantumHasWeight}}}
\expandafter\def\csname proof@DistanceMinima.sourceQuantumDistance\endcsname{\href{https://github.com/alectionary/vibe-math/blob/3ae21c0719cedbd928f4f7882f9b9e266567605e/formalization/distance-preservation/truncated12/DistanceMinima.lean\#L70-L71}{\nolinkurl{TruncatedDistance.sourceQuantumDistance}}}
\expandafter\def\csname proof@DistanceMinima.source_quantum_minimum_lower_iff\endcsname{\href{https://github.com/alectionary/vibe-math/blob/3ae21c0719cedbd928f4f7882f9b9e266567605e/formalization/distance-preservation/truncated12/DistanceMinima.lean\#L88-L103}{\nolinkurl{TruncatedDistance.source_quantum_minimum_lower_iff}}}
\expandafter\def\csname proof@DistanceTwelve.source_distance_bound_twelve\endcsname{\href{https://github.com/alectionary/vibe-math/blob/3ae21c0719cedbd928f4f7882f9b9e266567605e/formalization/distance-preservation/truncated12/DistanceTwelve.lean\#L9-L18}{\nolinkurl{TruncatedDistance.source_distance_bound_twelve}}}
\expandafter\def\csname proof@CriterionScope.no_full_row_of_no_triple\endcsname{\href{https://github.com/alectionary/vibe-math/blob/3ae21c0719cedbd928f4f7882f9b9e266567605e/formalization/distance-preservation/truncated12/CriterionScope.lean\#L12-L17}{\nolinkurl{TruncatedDistance.no_full_row_of_no_triple}}}

\lstset{basicstyle=\footnotesize\ttfamily,columns=fullflexible,keepspaces=true,
 frame=single,framerule=.3pt,breaklines=true,showstringspaces=false,
 escapeinside={(*@}{@*)},aboveskip=8pt,belowskip=10pt}
\title{A truncated distance bound for cyclic lifted-product quantum codes\\
Small-support geometry and checked integer certificates}
\author{}
\date{25 October 2026}
\begin{document}
\maketitle
\begin{abstract}
We prove a sufficient condition for the ordinary Pauli distance of a symmetric
cyclic lifted-product code to satisfy $d_Q\geq\min(d_C,12)$.
The base is a normalized, fully populated $3\times4$ monomial matrix over
$\F[x]/(x^L-1)$. The conditions are classical Tanner girth at least eight,
absence of a full column partition, exclusion of three equal values of each
three-row pairing form, and transpose-code distance at least $d_C$.
The proof works for arbitrary cyclic lift length, including composite lengths.
A source-kernel argument reduces every mixed support of weight at most eleven
to three block-count profiles, of weights seven, ten and eleven.
Actual support encodings and 183,312 checked integer certificates eliminate
the non-stabilizer possibilities. Six exceptional certificates instead force
a classical word of weight at most ten, which suffices for the distance bound.
Lean checks the complete implication for the actual expanded binary matrices,
both full stabilizer spans, both CSS sectors and Pauli union weight, together
with the literal distance minima and the transpose infinity convention.
The result is a bounded sufficient condition motivated by Raveendran,
Declercq and Vasi\'c's question about $d_Q=d_C$; it does not settle that broader
question.
\end{abstract}

\section{The theorem and the source question}\label{sec:answer}
Let $L\geq2$ and $G=\Z/L\Z$. A matrix
$E:\{0,1,2\}\times\{0,1,2,3\}\to G$ specifies the monomial base
$B_{ic}=x^{E_{ic}}$. Every entry is present: exponent zero means an identity
circulant. We normalize
\begin{equation}\label{eq:normal}
 E_{0c}=E_{i0}=0.
\end{equation}
The associated symmetric lifted product has $25L$ physical qubits and
binary checks obtained from
\begin{equation}\label{eq:source}
 C=\operatorname{expand}(B),\qquad
 \HX=\operatorname{expand}[B\otimes I_4\mid I_3\otimes B^*],\qquad
 \HZ=\operatorname{expand}[I_4\otimes B\mid B^*\otimes I_3].
\end{equation}
Here star transposes and sends $x$ to $x^{-1}$. Section~\ref{sec:model}
fixes the coefficient convention and the full logical quotient.

We use two explicit entry tests. The \emph{column test} excludes any partition
of all four distinct columns $a,b,c,d$ for which
\begin{equation}\label{eq:column}
 E_{ia}+E_{ib}=E_{ic}+E_{id}\quad\hbox{for every row }i.
\end{equation}
For distinct rows $i,j,k$, put
\begin{equation}\label{eq:rowform}
 F_{ijk}(c)=E_{ic}+E_{jc}-2E_{kc}.
\end{equation}
The \emph{triple-row test} says that no three distinct columns have the same
value of $F_{ijk}$. In particular, it excludes a full-row RCPC, since four
zero values would contain three equal values. These tests are stated directly
in the exponent group; their exact relation to the source conventions is
recorded in Appendix~\ref{app:scope}.

\begin{theorem}[Truncated source distance]\label{thm:main}
For the construction~\eqref{eq:source}, assume normalization~\eqref{eq:normal},
classical Tanner girth at least eight, the column and triple-row tests, and
$d_{C^{\mathsf T}}\geq d_C$. Then
\begin{equation}\label{eq:main}
                         d_Q\geq\min(d_C,12).
\end{equation}
The transpose distance is $+\infty$ when $\ker C^{\mathsf T}=0$.
Quantum distance is measured with both entire stabilizer row spans and
ordinary Pauli union weight. The implication holds for arbitrary cyclic
modulus; neither oddness nor primality is assumed.
\end{theorem}

The complete distance statement is
\lean{DistanceTwelve}{source_distance_bound_twelve}.
Its general predicate form,
\lean{ThroughEleven}{source_pauli_bound_twelve}, assumes that every nonzero
classical and transpose kernel word has weight at least $D$, and concludes
the Pauli bound $\min(D,12)$. The proof of Theorem~\ref{thm:main} is completed
in Section~\ref{sec:assembly}. Appendix~\ref{app:coverage} maps every argument
below to the checked declarations; the following Lean appendices print
complete central proof bodies and link every imported research module.

\subsection{What is answered}
Raveendran, Declercq and Vasi\'c~\cite[Eq.~(9), Definitions 2--5 and
Section VI]{RDV} ask for sufficient base constraints ensuring equality of the
quantum and underlying classical distances. Their partition constraints
motivate examining logical supports near the stabilizer-generator scale.
For a $3\times4$ base a generator has weight seven, whereas a mixed logical
can involve more coordinates. A previously established counterfamily
\cite{Counterfamily} has classical girth eight, no full RCPC and
$d_Q\leq11<12\leq d_C$, even when its transpose distance is at least $d_C$.
This motivates controlling partial row pairing as well as full partitions.

The triple-row condition and the truncated target are the bounded criterion
studied here, rather than a numbered conjecture attributed to the source
paper. The distinction matters: a lower bound cut off at twelve leaves all
higher-weight logical supports untreated and supplies no upper bound on $d_Q$.

\begin{longtable}{@{}P{.25\textwidth}P{.32\textwidth}P{.37\textwidth}@{}}
\toprule Question & Source and locator & Answer\\\midrule\endhead
Which base constraints ensure $d_Q=d_C$? & Raveendran--Declercq--Vasi\'c,
arXiv:2503.07567v1, 10 March 2025, Section VI. & Remains open in general here.\\[4pt]
Does the strengthened normalized $3\times4$ criterion imply
$d_Q\geq\min(d_C,12)$? & Project-proposed bounded follow-up to that question;
construction Eq.~(9), entry conventions Definitions 2--5. & Yes:
Theorem~\ref{thm:main}, including actual source matrices and complete
stabilizer quotients.\\[4pt]
Can the finite algebraic obstructions be interpreted for composite lifts? &
The proof obligation for this criterion. & Yes: integer identities are
interpreted without division, and the six torsion-five alternatives are
classical low-weight words (Lemma~\ref{lem:torsion}).\\
\bottomrule\end{longtable}

\subsection{Comparison and reference chronology}
Pairing selected coordinates at checks, passing to a colored triangle-free
normal graph and using exact integer consistency certificates are related
prior methods in Okada and Kasai~\cite[Sections 3.4 and 3.7]{OK}.
Their inspected v2 proves distance eighteen for one $[[34542,23032,18]]$
code and an affine-lift ceiling for a different fixed finite-field base at
prime lift degree greater than nineteen. Those statements do not establish
the normalized $3\times4$ cyclic criterion here. The contribution of the
present proof is its source-complete reduction and certificate interpretation
for this family at every cyclic modulus, with the full logical quotient.

Lin and Pryadko~\cite[Section IV.6, Statement 12]{LP} give distance bounds for
two-block group-algebra codes involving a support-intersection subgroup.
Their two-element construction and bound do not imply this matrix-base
criterion. Lee~\cite{Lee} supplies exact code-specific distance searches and
equivalence reductions for small two-block codes; those computations are
useful background, rather than premises of the universal implication here.
The tree counting behind paired checks and Mantel's triangle-free edge bound
are classical ingredients; related stopping-set tree bounds are discussed by
Kelley et al.~\cite{Kelley}. The actual matrix and quotient translations are
included in the checked proof.

The central chronology is shown below; Appendix~\ref{app:sources} specifies
coverage and date precision. Calendar ages are not durations of an unsolved
problem or measurements of proof effort.
\begin{center}\small
\begin{tabular}{@{}P{.25\textwidth}P{.20\textwidth}P{.23\textwidth}P{.23\textwidth}@{}}
\toprule Source & Earliest date & Inspected version & Age to 25 Oct. 2026\\\midrule
Raveendran et al.~\cite{RDV} & 10 Mar. 2025 & v1, 10 Mar. 2025 & 594 days\\
Lin--Pryadko~\cite{LP} & 28 June 2023 & v1, same date & 1,215 days\\
Okada--Kasai~\cite{OK} & 25 June 2026 & v2, 5 Sept. 2026 & 122 days; 50 from v2\\
Lee~\cite{Lee} & 9 Oct. 2026 & v2, 12 Oct. 2026 & 16 days; 13 from v2\\
Kelley et al.~\cite{Kelley} & 2004 (year only) & ISIT 2004 & About 22 years\\
\bottomrule\end{tabular}
\end{center}

\section{Source matrices, kernel words and the full quotient}\label{sec:model}
A classical variable is $(c,s)\in\{0,1,2,3\}\times G$ and a classical check is
$(i,t)\in\{0,1,2\}\times G$, with coefficient one precisely when
$t=s+E_{ic}$. Left product coordinates are $U(j,k,s)$, with $j,k\in\{0,1,2,3\}$;
right coordinates are $V(\ell,i,s)$, with $\ell,i\in\{0,1,2\}$. The actual
$Z$-check equation at $(j,i,t)$ is
\begin{equation}\label{eq:syndrome}
 \sum_k U(j,k,t-E_{ik})+\sum_\ell V(\ell,i,t+E_{\ell j})=0\quad\hbox{in }\F.
\end{equation}
A left coordinate meets three checks and a right coordinate meets four.
Both check matrices have $12L$ rows and $25L$ columns.

For binary vectors $x,z$ on these physical coordinates, let
\begin{equation}\label{eq:pauli}
 \HZ x=0,\qquad \HX z=0,\qquad
 (x,z)\notin\row(\HX)\times\row(\HZ).
\end{equation}
Their Pauli weight is $|\supp(x)\cup\supp(z)|$. Every linear combination of
every stabilizer row is allowed in~\eqref{eq:pauli}; no restricted collection
of generators is substituted. The quantum distance is the minimum of these
weights. Equivalently, one may take the extended-natural infimum and assign
$+\infty$ if there is no nontrivial pair. The lower bound applies to that
convention as well as the usual positive-logical-dimension regime.

The classical distance is the least nonzero weight in $\ker C$, and the
transpose distance uses $\ker C^{\mathsf T}$. A pure-left source kernel word
splits into classical kernel words, one for each $j$; a pure-right word
splits into transpose kernel words, one for each $i$. Thus a nonzero pure
word cannot have weight below $D$ if both classical lower predicates hold.
The checked pure-sector statements are used in Lemma~\ref{lem:global}.

\section{Small supports have only three mixed profiles}\label{sec:geometry}
\begin{lemma}[Actual incidence and paired checks]\label{lem:incidence}
Classical no-four/no-six implies no-four/no-six for the actual $Z$-check
incidence of~\eqref{eq:source}. In a source kernel word of weight at most
eleven, every check meets its support in zero or two coordinates.
\end{lemma}
\begin{proof}
The coefficient incidence is exactly the relation in~\eqref{eq:syndrome}.
Two coordinates of the same product block sharing two checks give a classical
four-cycle; a mixed pair has only one check type in common. A six-cycle
contained in one block gives a classical six-cycle. In a mixed six-cycle,
the lone coordinate from one block would have to use the same check type
twice, contradicting distinctness of the cycle checks.

Kernel parity makes each selected check degree even. If a check met at least
four selected coordinates, choose four. Each has at least two other checks;
each of those needs another selected neighbor. Four-cycle exclusion makes
the eight other checks distinct, and four-/six-cycle exclusion makes their
eight required support neighbors distinct from one another and from the
four original coordinates. The support would have at least twelve elements.
The actual source-degree and short-cycle premises of this tree argument are
proved in \texttt{ActualIncidence}; its conclusion is
\lean{ActualIncidence}{source_kernel_zero_or_two}.
\end{proof}

At each occupied check, join its two selected coordinates. This normal graph
is simple by four-cycle exclusion and triangle-free by six-cycle exclusion.
Its left and right vertex degrees are three and four. Set
\begin{equation}\label{eq:counts}
 a_j=|\supp U\cap\{\hbox{first index }j\}|,\qquad
 b_i=|\supp V\cap\{\hbox{second index }i\}|.
\end{equation}
\begin{lemma}[Complete mixed count classification]\label{lem:profiles}
For a source kernel word of weight at most eleven with both blocks nonempty,
the separately sorted counts are exactly one of
\begin{equation}\label{eq:profiles}
 \begin{array}{c|c|c}
 \hbox{weight}&a&b\\\hline
 7&(1,1,1,1)&(1,1,1)\\
 10&(2,2,2,0)&(2,2,0)\\
 11&(5,1,1,1)&(1,1,1).
 \end{array}
\end{equation}
\end{lemma}
\begin{proof}
Sum~\eqref{eq:syndrome} over the fiber variable. For every $j,i$ it gives
$a_j+b_i=0\pmod2$, so all seven counts have one common parity.
Let $m_{ji}$ count internal left-left normal edges in bucket $(j,i)$ and
$c_{ji}$ count cross edges there. Socket counting gives
$a_j=2m_{ji}+c_{ji}$ and $c_{ji}\leq b_i$. Distinct internal edges of a fixed
left group form a triangle-free graph on $a_j$ vertices. Mantel's bound gives
\begin{equation}\label{eq:mantel}
 \sum_i\frac{\max(a_j-b_i,0)}2\leq\left\lfloor\frac{a_j^2}4\right\rfloor,
 \qquad
 \sum_j\frac{\max(b_i-a_j,0)}2\leq\left\lfloor\frac{b_i^2}4\right\rfloor.
\end{equation}
All differences being divided by two are even. The analogous right-group
argument proves the second inequality.
There are forty-four sorted common-parity profiles with both blocks nonempty
and total at most eleven. An exhaustive bounded-list enumeration proves
completeness, and checking~\eqref{eq:mantel} on those profiles leaves exactly
\eqref{eq:profiles}. The enumeration and inequalities are both checked from
the actual source kernel in
\lean{SourceProfiles}{source_mixed_count_profiles}; the independent socket
classification agrees. The enumeration is over integer counts, not sampled
bases or lift lengths.
\end{proof}

\section{Canonical source supports and finite encodings}\label{sec:encoding}
\begin{lemma}[Transport preserves the complete problem]\label{lem:transport}
Row/column permutation followed by re-normalization preserves the entry tests,
classical and transpose lower predicates, source kernel, full stabilizer
membership, weight and group counts. Hence the missing weight-ten groups
may be placed at zero and the exceptional weight-eleven group at index three.
\end{lemma}
\begin{proof}
Permutations relabel both tensor factors. For a gauge change
$E'_{ij}=E_{ij}+r_i+c_j$, use the fiber bijections
\begin{equation}\label{eq:gauge}
 U'(j,k,s)=U(j,k,s+c_j+c_k),\qquad
 V'(\ell,i,s)=V(\ell,i,s-r_\ell-r_i).
\end{equation}
Substitution checks the kernel equation; stabilizer coefficients transform
by $w'(i,j,s)=w(i,j,s-r_i+c_j)$. These are coordinate bijections, so both
full spans and weight are preserved. Classical and transpose kernels have
the corresponding coordinate transports.

Re-normalizing a permuted base subtracts one column-independent constant
from each row form, preserving equality of three values. For the column
partition form $A_i=E_{ia}+E_{ib}-E_{ic}-E_{id}$ it replaces $A_i$ by
$A_i-A_r$. The original zero row gives $A_0=0$, so simultaneous vanishing
is preserved. Girth is preserved by the source incidence bijections.
Finally the count classification supplies the unique missing or exceptional
group and the values of every other group. The complete transports are in
\texttt{ProfileTransport}; the marker choices are in \texttt{CountCanonical}.
Thus canonical placement is derived, rather than an additional hypothesis.
\end{proof}

\begin{lemma}[Singleton labels and the constant-row branch]\label{lem:singleton}
The singleton groups in the weight-seven and weight-eleven profiles have
label maps that are either both constant or both injective. The constant
case is in the full stabilizer span when the word is below the classical
lower bound. The weight-seven profile is always a stabilizer under the
entry tests.
\end{lemma}
\begin{proof}
Write a singleton left coordinate as $(j,k_j,s_j)$ and a right coordinate
as $(\ell_i,i,t_i)$. Their actual cross checks give
\begin{equation}\label{eq:cross}
 s_j+E_{i,k_j}=t_i-E_{\ell_i,j}.
\end{equation}
If two $k_j$ agree, subtraction forces a column difference to be independent
of $i$; classical four-cycle exclusion forces all $\ell_i$ to agree.
The converse is symmetric. Hence the maps are both constant or both injective.

For constants $K,H$, equation~\eqref{eq:cross} supplies a common $z$ with
$s_j=z-E_{Hj}$ and $t_i=z+E_{iK}$. These are precisely the coordinates of
the complete $X$-stabilizer row indexed by $(H,K,z)$. At weight seven this
is the whole word. At weight eleven, adding that row cancels the three
singleton left coordinates and all three right coordinates; the remainder
is pure-left of weight at most six. Below the classical lower bound it
vanishes, proving full stabilizer membership. The checked source statement
includes the right coordinates and all stabilizer coefficients.

In the injective weight-seven case, the label maps are permutations of four
columns and three rows. The resulting $24\cdot6=144$ patterns are all
certified forbidden by the entry/girth tests through the integer argument
of Section~\ref{sec:certificates}. Actual source realizations are constructed
in \texttt{SourceSeven}; none is postulated as a new premise.
\end{proof}

\begin{lemma}[Complete weight-ten encoding]\label{lem:ten}
Every actual weight-ten kernel support with counts in~\eqref{eq:profiles},
after the transport of Lemma~\ref{lem:transport}, realizes an index in a
catalog of size $124,416$.
\end{lemma}
\begin{proof}
At the missing right row zero, each remaining two-coordinate left group
must be internally paired. Its two fibers agree by normalization, and its
two column labels differ. The same two coordinates cannot be internally
paired at another row, since that would give a classical four-cycle.
Similarly each right pair is internally matched at missing left column zero,
with equal fibers and different row labels. Every other bucket is a
bijection between two-element sets and has exactly two matching choices.

The three increasingly ordered column pairs are selected from six possibilities each,
the two increasingly ordered row pairs from three possibilities each, and the six
cross-bucket matchings from two possibilities each. Thus the finite decoder
has $6^3 3^2 2^6=124,416$ indices. The geometry is extracted from actual
supports; rank bounds, decoder identities and every source incidence are
proved before applying the catalog. The complete interface is
\lean{TenEncoding}{actual_canonical_ten_support_encodes}.
\end{proof}

\begin{lemma}[Complete weight-eleven encoding]\label{lem:eleven}
For a weight-eleven support in~\eqref{eq:profiles}, the constant-label branch
is covered by Lemma~\ref{lem:singleton}. Every injective-label support realizes
an index in a catalog of size $58,752$.
\end{lemma}
\begin{proof}
The exceptional left group has five vertices and six internal normal edges:
its fifteen sockets have exactly three cross partners. A triangle-free
five-vertex graph with six edges is $K_{2,3}$. Its two centers use all three
row colors internally; each leaf has one missing color and one cross partner.
The missing colors are distinct because the right singleton in any row has
only one socket at the exceptional first-index group. Index the leaves by
those missing colors. There are two compatible row-color orientations.
Adjacent internal vertices cannot have the same second column label, since
their common source check would then force identical coordinates.

The complete finite label table has 204 admissible center/leaf tuples.
There are 24 injective singleton-column maps, six row permutations and two
orientations, giving $24\cdot6\cdot204\cdot2=58,752$ decoder indices.
The actual $K_{2,3}$ numbering, colors, missing-row permutation, singleton
partners and eighteen source incidences are extracted in the source proof.
\lean{SourceEleven}{source_canonical_eleven_stabilizer} assembles this encoding
with the complete constant-label branch and the certified exclusions.
\end{proof}

\section{Integer certificates for every encoded pattern}\label{sec:certificates}
\begin{lemma}[Universal interpretation of the certificates]\label{lem:cert}
A checked integer row identity for a catalog pattern implies its specified
forbidden equation in $G$, whenever that pattern has an actual source
realization. This implication requires no division and holds at every cyclic
modulus.
\end{lemma}
\begin{proof}
Let $y$ consist of the six normalized exponents and the $w$ vertex fibers.
Every matching edge gives a homogeneous integer equation in these variables,
with signs fixed by~\eqref{eq:syndrome}. The full matching matrix $M$ uses
eighteen row slots, padding unused slots when necessary. A certificate for
a target row $f$ is an integer vector $a$ satisfying
$f_j=\sum_r a_rM_{rj}$ for every variable $j$. In every additive group,
\[
 \sum_j f_j y_j=\sum_r a_r\Bigl(\sum_j M_{rj}y_j\Bigr)=0.
\]
All fiber columns are included in the checked identity. Certificates involving
several targets provide each identity separately.

The checker verifies these integer equalities, and
\lean{KernelCatalog}{checked_fast_equations} interprets its Boolean success
through \lean{IntegerSoundness}{integer_certificate_sound}.
An actual realization supplies the matching equations. The sixty-four target
outcome types are interpreted as short cycles, a prohibited entry pattern,
small cyclic torsion, or the classical alternative in Lemma~\ref{lem:torsion}.
This final source exclusion is
\lean{PatternRealizations}{certified_pattern_has_no_source_realization}.
\end{proof}

\begin{lemma}[Cyclic torsion and the six classical alternatives]\label{lem:torsion}
The small row/column torsion outcomes contradict classical no-four.
Each of the six remaining weight-ten alternatives instead forces a nonzero
classical word of weight at most ten, contradicting $10<d_C$.
\end{lemma}
\begin{proof}
For distinct rows, the four differences across columns are distinct under
no-four; for distinct columns, the three row differences are distinct.
For positive $n$, the $n$-torsion subgroup of $G$ has at most $n$ elements. Consequently four
distinct row differences cannot all be killed by $1\leq n\leq3$, and three distinct
column differences cannot all be killed by two.

Each exceptional weight-ten pattern forces $5E_{1c}=5E_{2c}=0$ for some
$c\ne0$. Set $K=\{s\in G:5s=0\}$ and take the binary indicator of
$\{0,c\}\times K$. It is nonzero and has weight $2|K|\leq10$.
In each classical check its two contributions cancel: $E_{i0}=0$, and
translation by $E_{ic}\in K$ preserves $K$. This is an actual $C$-kernel word,
proved by \lean{TorsionClassical}{source_classical_word_le10}.
The alternative contradicts a proposed weight-ten support strictly below
$d_C$; it is not a qualifying counterexample. In particular the argument
never divides by five, including when $5\mid L$.
\end{proof}

\begin{proposition}[Complete quantified catalog coverage]\label{prop:coverage}
Every decoder index at weights seven, ten and eleven has a kernel-checked
certificate. The catalog sizes are respectively $144$, $124,416$ and $58,752$.
\end{proposition}
\begin{proof}
Each closed Lean case proves that the full integer-identity checker returns
true for its explicit certificate. A balanced interval proof assembles the
closed cases into $\forall p:\operatorname{Fin}(N),\exists c,\mathrm{check}(p,c)$.
The checked append lemma composes adjacent intervals and covers both branches
of the index split. The resulting uniform statements are
\lean{SourceSeven}{seven_certificates_complete},
\lean{Coverage10}{complete_coverage10} and
\lean{Coverage11}{complete_coverage11}.
An independent standard-library verifier reconstructs every matching matrix
and integer identity; it agrees on all $183,312$ indices, but is not a premise
of these Lean theorems.
\end{proof}

\begin{center}\small
\begin{tabular}{@{}lrrr@{}}
\toprule Certificate outcome & Weight 7 & Weight 10 & Weight 11\\\midrule
Classical four-cycle &109&124,129&57,282\\
Classical six-cycle &16&262&1,008\\
Forbidden three-column row form &9&9&342\\
Full column partition &3&3&66\\
Small cyclic row torsion &1&1&0\\
Column two-torsion &6&6&54\\
Classical word of weight at most ten &0&6&0\\\midrule
Total &144&124,416&58,752\\\bottomrule
\end{tabular}
\end{center}
These are certificate classifications of the finite decoder indices. They
are not counts of inequivalent codes, lift lengths or logical operators.
Discovery used integer algebra; checking uses the full source matching rows,
including fiber variables. No finite case remains unclassified.

\section{The full source bound and ordinary Pauli distance}\label{sec:assembly}
\begin{lemma}[Global full-span reduction]\label{lem:global}
Assume the source entry/girth tests and classical/transpose lower predicates
at $D$. Every source $Z$-kernel word of weight below $\min(D,12)$ lies in the
entire row span of $\HX$.
\end{lemma}
\begin{proof}
A pure-block word below $D$ is zero by the actual classical or transpose
kernel bound. Otherwise Lemma~\ref{lem:profiles} leaves only weights seven,
ten and eleven. Lemma~\ref{lem:transport} supplies canonical placement and
preserves every hypothesis and the full quotient.
The weight-seven case is a stabilizer by Lemma~\ref{lem:singleton}.
For weight ten, Lemma~\ref{lem:ten} gives an actual realization of a certified
pattern. Its outcome contradicts a source test, or forces $d_C\leq10$ while
the word has weight ten below $D$. For weight eleven the constant case is a
stabilizer; the injective case is excluded by Lemma~\ref{lem:eleven} and the
complete catalog. Transport the resulting full-span statement back to the
original source coordinates. The case split is printed in the Lean appendix
as \lean{TwelveConditional}{source_kernel_below_twelve_of_coverage};
\lean{ThroughEleven}{source_kernel_below_twelve_is_stabilizer} discharges its
finite-coverage premise with Proposition~\ref{prop:coverage}.
\end{proof}

\begin{lemma}[Both sectors and union weight]\label{lem:pauli}
An $X$-logical lower bound $b$ for this source construction gives the same
$Z$-logical lower bound and the ordinary Pauli lower bound $b$.
\end{lemma}
\begin{proof}
Swap the two finite indices within each tensor-coordinate block.
This involutive coordinate permutation exchanges the actual $X$/$Z$ kernel
equations and full spans and preserves support size. A nontrivial Pauli pair
has a component outside its corresponding full span. That component is a
pure logical, and its weight is at most the union weight. Apply the sector
bound and support inclusion. The complete union-weight proof is
\lean{PauliBridge}{pauli_bound_of_x_bound}.
\end{proof}

\begin{proposition}[Literal distance minima]\label{prop:minima}
The classical minimum is attained and satisfies its universal lower predicate.
The transpose minimum is infinity for a zero transpose kernel, and a comparison
$D\leq d_{C^{\mathsf T}}$ implies the corresponding lower predicate.
For the full Pauli quotient, $b\leq d_Q$ is equivalent to saying that every
nontrivial pair has union weight at least $b$.
\end{proposition}
\begin{proof}
The actual $3L\times4L$ classical matrix has a nonzero kernel word, as proved
in \texttt{GeneralLPPhysical}; define its distance by the least attained
natural weight. Leastness gives the universal predicate. Define transpose
and Pauli distances as extended-natural infima of their actual nonzero-word
or nontrivial-pair weights. The infimum of the empty set is infinity.
An infimum is at most each member, and a lower bound on every member is at
most the infimum. These properties prove the required implications and the
Pauli equivalence. All three definitions and the bridges are checked in
\texttt{DistanceMinima}; the public theorem retains the complete quotient.
\end{proof}

\begin{proof}[Proof of Theorem~\ref{thm:main}]
Set $D=d_C$. Proposition~\ref{prop:minima} supplies the classical lower
predicate and obtains the transpose predicate from the assumed comparison.
Lemma~\ref{lem:global} excludes every $X$ logical below $\min(D,12)$;
Lemma~\ref{lem:pauli} gives the ordinary Pauli bound. Apply the minimum
characterization in Proposition~\ref{prop:minima}. This is exactly the
proof body of \lean{DistanceTwelve}{source_distance_bound_twelve}.
\end{proof}

\section{Implications and remaining question}\label{sec:limits}
The theorem controls the entire small-support range relevant to the
truncation: every hypothetical nontrivial word below $\min(d_C,12)$ is
excluded in the actual source model. The integer argument remains valid
at even and composite lift lengths; the torsion-five alternatives explain
why a certificate need not be a pure entry obstruction to finish a distance
lower bound.

The conclusion is sufficient and one-sided. It does not characterize all
bases with preserved distance, give an exact $d_Q$, or settle the source's
unrestricted $d_Q=d_C$ design question. No decoder-performance conclusion
follows from it. Higher-weight supports need further arguments.
A concrete possible follow-up is the already-studied $L=13$ base with
classical distance fourteen and a saved quantum upper witness of weight
fourteen. Its computational exclusion through thirteen is separate evidence;
a faithful Lean distance-fourteen certificate would still need the concrete
hypotheses, upper witness and weight-twelve/thirteen exclusion. That finite
follow-up is not required for Theorem~\ref{thm:main}.

\appendix
\clearpage\section{Proof coverage and reproducibility}\label{app:coverage}
The checked toolchain is Lean 4.36.0-rc1, commit
\nolinkurl{c9afcddf514a6772bd4d649f7cfc04eff0b9f15e}, with Mathlib commit
\nolinkurl{5f7b4198d22a9c82e1c698a379334fe128095c82}.
The immutable proof snapshot is
\nolinkurl{3ae21c0719cedbd928f4f7882f9b9e266567605e}.
The main source-distance theorem and source-scope audit have 269 research
imports, all admitted against matching source and log hashes. Their reported
axioms are contained in \texttt{propext}, \texttt{Classical.choice} and
\texttt{Quot.sound}. No native decision oracle or unproved catalog premise
is admitted. The dependency cache is pinned; Mathlib itself is not presented
as newly rebuilt here.

The table maps every theorem/proof in the mathematical exposition to its
faithful formal counterpart, including the support constructions, finite
coverage and literal minima. \texttt{ReportAudit.lean} checks the named
declarations and prints their axiom closures. The following appendices print
complete central declarations extracted from the immutable checked sources.
The large finite case files are linked in full in Appendix~\ref{app:links}.
\begingroup\footnotesize\begin{longtable}{@{}P{.30\textwidth}P{.65\textwidth}@{}}
\toprule Manuscript proof & Checked Lean declarations\\\midrule\endhead
\ref{thm:main}: Literal source-distance bound & \lean{DistanceTwelve}{source_distance_bound_twelve}\\[5pt]
\ref{lem:incidence}: Actual coefficients, short cycles and paired checks & \lean{ActualIncidence}{classical_incidence}; \lean{ActualIncidence}{quantum_incidence}; \lean{ActualIncidence}{quantum_no_four}; \lean{ActualIncidence}{quantum_no_six}; \lean{ActualIncidence}{source_kernel_zero_or_two}\\[5pt]
\ref{lem:profiles}: Parity, local Mantel bounds and exhaustive profiles & \lean{BlockParity}{common_block_parity}; \lean{BlockCounts}{common_count_parity}; \lean{LocalMantel}{source_left_local_mantel}; \lean{LocalMantel}{source_right_local_mantel}; \lean{CountPatterns}{count_cases_length}; \lean{CountPatterns}{classified_of_mem_countCases}; \lean{SourceProfiles}{source_mixed_count_profiles}\\[5pt]
\ref{lem:transport}: Full source transport and canonical placement & \lean{ProfileTransport}{normPerm_criterion}; \lean{ProfileTransport}{normPerm_kernel}; \lean{ProfileTransport}{normPerm_span}; \lean{ProfileTransport}{normPerm_weight}; \lean{CountCanonical}{ten_left_marker}; \lean{CountCanonical}{ten_right_marker}; \lean{CountCanonical}{eleven_left_marker}\\[5pt]
\ref{lem:singleton}: Constant or injective singleton labels and full-row cancellation & \lean{SupportSingletons}{actual_singleton_label_alternative}; \lean{ConstantRows}{constant_singleton_rows}; \lean{RowCancellation}{single_exception_cancellation}; \lean{SourceSeven}{source_seven_singleton_profile_stabilizer}\\[5pt]
\ref{lem:ten}: Actual paired weight-ten source encoding & \lean{TenGeometry}{source_ten_geometry}; \lean{TenEncoding}{actual_canonical_ten_support_encodes}; \lean{TenIndex}{tenRank_bound}\\[5pt]
\ref{lem:eleven}: Actual colored K2,3 and weight-eleven source encoding & \lean{ExceptionalGraph}{exceptional_group_is_k23}; \lean{ElevenEncoding}{eleven_geometry_encodes}; \lean{SourceEleven}{source_canonical_eleven_stabilizer}\\[5pt]
\ref{lem:cert}: Full integer identities, universal interpretation and source obstruction & \lean{IntegerSoundness}{integer_certificate_sound}; \lean{KernelCatalog}{checked_fast_equations}; \lean{PatternRealizations}{certified_pattern_has_no_source_realization}; \lean{TargetObstructions}{source_target_obstruction}\\[5pt]
\ref{lem:torsion}: Classical alternatives without division & \lean{TorsionClassical}{torsion_card_le}; \lean{TorsionClassical}{torsionWord_source_kernel}; \lean{TorsionClassical}{source_classical_word_le10}; \lean{IntegerSoundness}{row_small_torsion_impossible}; \lean{IntegerSoundness}{column_two_torsion_impossible}\\[5pt]
\ref{prop:coverage}: Quantified certificate coverage & \lean{SourceSeven}{seven_certificates_complete}; \lean{Coverage10}{complete_coverage10}; \lean{Coverage11}{complete_coverage11}; \lean{CoverageCore}{covered_append}\\[5pt]
\ref{lem:global}: Pure sectors and complete full-span reduction & \lean{PureDistance}{pure_left_zero_below_classical}; \lean{PureDistance}{pure_right_zero_below_transpose}; \lean{TwelveConditional}{source_kernel_below_twelve_of_coverage}; \lean{ThroughEleven}{source_kernel_below_twelve_is_stabilizer}\\[5pt]
\ref{lem:pauli}: Actual sector symmetry and ordinary union weight & \lean{GeneralSymmetry}{full_span_swap}; \lean{GeneralSymmetry}{kernel_swap}; \lean{GeneralSymmetry}{swap_weight}; \lean{PauliBridge}{pauli_bound_of_x_bound}; \lean{ThroughEleven}{source_pauli_bound_twelve}\\[5pt]
\ref{prop:minima}: Actual classical/transpose/quantum minima & \lean{DistanceMinima}{source_classical_minimum_realized}; \lean{DistanceMinima}{source_classical_minimum_lower}; \lean{DistanceMinima}{source_transpose_zero_kernel_infinity}; \lean{DistanceMinima}{source_quantum_minimum_lower_iff}\\[5pt]
\bottomrule\end{longtable}
\endgroup

The proof dependencies can be read in the following order:
\begin{center}
\begin{tabular}{c}
Actual source matrices and incidence\\$\downarrow$\\
Paired checks, parity and local triangle-free edge bounds\\$\downarrow$\\
Three count profiles; actual weight-seven/ten/eleven encodings\\$\downarrow$\\
Full integer identities and quantified finite coverage\\$\downarrow$\\
Permutation/normalization transport and full-span reduction\\$\downarrow$\\
Both sectors, union weight and literal distance minima
\end{tabular}
\end{center}

From the investigation directory, run the independent verifier and rebuild
with pinned dependencies in a fresh object directory:
\begin{lstlisting}
python3 verify_catalog.py
python3 replay_round.py --mathlib /path/to/pinned/mathlib4 \
  --lean-bin /path/to/pinned/lean/bin --build-dir /path/to/fresh/objects \
  DistanceTwelve CriterionScope
python3 verify_round.py --round-number 2 --final
\end{lstlisting}
The archive preserves relative source relationships and all required research
imports. Generators, encoded certificates, interval plans and build records
are included. Discovery outputs and prior unsuccessful classifications remain
in the research repository; they are not proof premises. The standalone
report build and integrity checker are described in the bundled
\texttt{README.txt}. Investigation effort and report preparation timestamps
are recorded outside the mathematical exposition in the provenance records.

\clearpage\section{Source scope and reference ledger}\label{app:scope}
The source construction is Raveendran--Declercq--Vasi\'c v1,
equation (9). The column condition is its full-column RCPC exclusion at
four columns (Definition 2). At three rows, full-row RCPC has the row form
\eqref{eq:rowform} equal to zero at every column (Definition 5).
The stronger triple-row test excludes that case;
\lean{CriterionScope}{no_full_row_of_no_triple} checks the implication.
Thus omission of a separate full-row field in the final Lean statement
removes only a redundant hypothesis. The formal theorem uses a positive
cyclic modulus; the displayed $L\geq2$ scope is contained in it.
The actual coefficient matrices, not just a support surrogate, determine all
kernel and span predicates.

\subsection{Chronology and bounded prior-art coverage}\label{app:sources}
The manuscript date is 25 October 2026. Exact calendar deltas in the table
near the introduction are differences to that date, using the earliest
arXiv posting and the inspected version date separately. The source histories
are linked in the bibliography and recorded in
\texttt{reference-chronology.json}. The 2004 stopping-set reference has only
year precision here, so its age is approximate. A source's age is not a
claim that its broad question was unsolved throughout that interval.

The source comparison reuses the completed full-text screens in
\href{https://github.com/alectionary/vibe-math/blob/3ae21c0719cedbd928f4f7882f9b9e266567605e/literature/truncated-distance12-source-screening.json}{the source ledger}.
It includes the original question, full-theorem comparisons and harvested
questions from encountered papers. The report metadata pass rechecks dates
of those same versions; it creates no new paper/question cohort.
Public GitHub and primary-source searches were bounded. Native X access was
unavailable; indexed discovery was not treated as authenticated coverage.
Uninspected leads remain named gaps in the ledger. These limits concern
source coverage, rather than a premise of the mathematical theorem.

\begin{thebibliography}{9}
\bibitem{RDV} N. Raveendran, D. Declercq and B. Vasi\'c,
\emph{On the Minimum Distances of Finite-Length Lifted Product Quantum LDPC Codes},
\href{https://arxiv.org/abs/2503.07567v1}{arXiv:2503.07567v1}, 10 March 2025.
Construction Eq.~(9), Definitions 2--5, Sections V-B and VI.
\bibitem{LP} H.-K. Lin and L. P. Pryadko,
\emph{Quantum two-block group algebra codes},
\href{https://arxiv.org/abs/2306.16400v1}{arXiv:2306.16400v1}, 28 June 2023.
Section IV.6, Statement 12.
\bibitem{OK} K. Okada and K. Kasai,
\emph{Rate-2/3 Girth-8 (3,18)-Regular Quantum LDPC Codes from Two-Branch
Finite-Field Bases and CPM Lifts},
\href{https://arxiv.org/abs/2606.27130v2}{arXiv:2606.27130v2},
first posted 25 June 2026; inspected v2, 5 September 2026.
\bibitem{Lee} G. Lee,
\emph{Bicycle Race: small two-block codes on the Pareto frontier},
\href{https://arxiv.org/abs/2610.13880v2}{arXiv:2610.13880v2},
first posted 9 October 2026; inspected v2, 12 October 2026.
\bibitem{Kelley} C. Kelley, D. Sridhara, J. Xu and J. Rosenthal,
\emph{Pseudocodeword Weights and Stopping Sets}, ISIT 2004.
\href{https://user.math.uzh.ch/kelley/isit2004_Chicago.pdf}{Author-hosted paper};
Theorem III.1 and its following note.
\bibitem{Counterfamily} \emph{Distance loss in girth-eight cyclic lifted-product
codes: An infinite counterfamily with checked Lean proofs}, repository report,
25 October 2026.
\href{https://github.com/alectionary/vibe-math/blob/3ae21c0719cedbd928f4f7882f9b9e266567605e/formalization/distance-preservation/paper/main.pdf}{Report and source links}.
\end{thebibliography}

\clearpage\section{Source matrices, full spans and exact distances}

The source interface fixes actual binary coefficients, every stabilizer row combination and the extended-natural minimum convention.

\subsection*{\texttt{GeneralLP}: HX}

\noindent\href{https://github.com/alectionary/vibe-math/blob/3ae21c0719cedbd928f4f7882f9b9e266567605e/formalization/distance-preservation/GeneralLP.lean#L11-L13}{Complete checked declaration at the immutable proof commit.}

\lstinputlisting{listings/GeneralLP-HX.tex}

\subsection*{\texttt{GeneralLP}: HZ}

\noindent\href{https://github.com/alectionary/vibe-math/blob/3ae21c0719cedbd928f4f7882f9b9e266567605e/formalization/distance-preservation/GeneralLP.lean#L15-L17}{Complete checked declaration at the immutable proof commit.}

\lstinputlisting{listings/GeneralLP-HZ.tex}

\subsection*{\texttt{GeneralLP}: InSpanX}

\noindent\href{https://github.com/alectionary/vibe-math/blob/3ae21c0719cedbd928f4f7882f9b9e266567605e/formalization/distance-preservation/GeneralLP.lean#L65-L67}{Complete checked declaration at the immutable proof commit.}

\lstinputlisting{listings/GeneralLP-InSpanX.tex}

\subsection*{\texttt{GeneralSymmetry}: fullSpanX}

\noindent\href{https://github.com/alectionary/vibe-math/blob/3ae21c0719cedbd928f4f7882f9b9e266567605e/formalization/distance-preservation/truncated12/GeneralSymmetry.lean#L22-L23}{Complete checked declaration at the immutable proof commit.}

\lstinputlisting{listings/GeneralSymmetry-fullSpanX.tex}

\subsection*{\texttt{GeneralSymmetry}: fullSpanZ}

\noindent\href{https://github.com/alectionary/vibe-math/blob/3ae21c0719cedbd928f4f7882f9b9e266567605e/formalization/distance-preservation/truncated12/GeneralSymmetry.lean#L24-L25}{Complete checked declaration at the immutable proof commit.}

\lstinputlisting{listings/GeneralSymmetry-fullSpanZ.tex}

\subsection*{\texttt{DistanceMinima}: sourceClassicalHasWeight}

\noindent\href{https://github.com/alectionary/vibe-math/blob/3ae21c0719cedbd928f4f7882f9b9e266567605e/formalization/distance-preservation/truncated12/DistanceMinima.lean#L9-L11}{Complete checked declaration at the immutable proof commit.}

\lstinputlisting{listings/DistanceMinima-sourceClassicalHasWeight.tex}

\subsection*{\texttt{DistanceMinima}: sourceClassicalDistance}

\noindent\href{https://github.com/alectionary/vibe-math/blob/3ae21c0719cedbd928f4f7882f9b9e266567605e/formalization/distance-preservation/truncated12/DistanceMinima.lean#L21-L22}{Complete checked declaration at the immutable proof commit.}

\lstinputlisting{listings/DistanceMinima-sourceClassicalDistance.tex}

\subsection*{\texttt{DistanceMinima}: sourceTransposeHasWeight}

\noindent\href{https://github.com/alectionary/vibe-math/blob/3ae21c0719cedbd928f4f7882f9b9e266567605e/formalization/distance-preservation/truncated12/DistanceMinima.lean#L33-L35}{Complete checked declaration at the immutable proof commit.}

\lstinputlisting{listings/DistanceMinima-sourceTransposeHasWeight.tex}

\subsection*{\texttt{DistanceMinima}: sourceTransposeDistance}

\noindent\href{https://github.com/alectionary/vibe-math/blob/3ae21c0719cedbd928f4f7882f9b9e266567605e/formalization/distance-preservation/truncated12/DistanceMinima.lean#L37-L38}{Complete checked declaration at the immutable proof commit.}

\lstinputlisting{listings/DistanceMinima-sourceTransposeDistance.tex}

\subsection*{\texttt{DistanceMinima}: sourceQuantumHasWeight}

\noindent\href{https://github.com/alectionary/vibe-math/blob/3ae21c0719cedbd928f4f7882f9b9e266567605e/formalization/distance-preservation/truncated12/DistanceMinima.lean#L60-L65}{Complete checked declaration at the immutable proof commit.}

\lstinputlisting{listings/DistanceMinima-sourceQuantumHasWeight.tex}

\subsection*{\texttt{DistanceMinima}: sourceQuantumDistance}

\noindent\href{https://github.com/alectionary/vibe-math/blob/3ae21c0719cedbd928f4f7882f9b9e266567605e/formalization/distance-preservation/truncated12/DistanceMinima.lean#L70-L71}{Complete checked declaration at the immutable proof commit.}

\lstinputlisting{listings/DistanceMinima-sourceQuantumDistance.tex}

\clearpage\section{Actual incidence and the three profiles}

The source kernel is the input of the count classification. The abstract normal graph and local edge bounds are proved in its imports, rather than assumed as new hypotheses.

\subsection*{\texttt{ActualIncidence}: source\_kernel\_zero\_or\_two}

\noindent\href{https://github.com/alectionary/vibe-math/blob/3ae21c0719cedbd928f4f7882f9b9e266567605e/formalization/distance-preservation/truncated12/ActualIncidence.lean#L127-L144}{Complete checked declaration at the immutable proof commit.}

\lstinputlisting{listings/ActualIncidence-source_kernel_zero_or_two.tex}

\subsection*{\texttt{SourceProfiles}: source\_mixed\_count\_profiles}

\noindent\href{https://github.com/alectionary/vibe-math/blob/3ae21c0719cedbd928f4f7882f9b9e266567605e/formalization/distance-preservation/truncated12/SourceProfiles.lean#L29-L79}{Complete checked declaration at the immutable proof commit.}

\lstinputlisting{listings/SourceProfiles-source_mixed_count_profiles.tex}

\clearpage\section{Constant rows and canonical support encodings}

The following complete declarations show the actual support-to-catalog interfaces, including the exceptional constant-label branch.

\subsection*{\texttt{ConstantRows}: constant\_singleton\_rows}

\noindent\href{https://github.com/alectionary/vibe-math/blob/3ae21c0719cedbd928f4f7882f9b9e266567605e/formalization/distance-preservation/truncated12/ConstantRows.lean#L25-L43}{Complete checked declaration at the immutable proof commit.}

\lstinputlisting{listings/ConstantRows-constant_singleton_rows.tex}

\subsection*{\texttt{SourceSeven}: source\_seven\_singleton\_profile\_stabilizer}

\noindent\href{https://github.com/alectionary/vibe-math/blob/3ae21c0719cedbd928f4f7882f9b9e266567605e/formalization/distance-preservation/truncated12/SourceSeven.lean#L258-L281}{Complete checked declaration at the immutable proof commit.}

\lstinputlisting{listings/SourceSeven-source_seven_singleton_profile_stabilizer.tex}

\subsection*{\texttt{TenGeometry}: source\_ten\_geometry}

\noindent\href{https://github.com/alectionary/vibe-math/blob/3ae21c0719cedbd928f4f7882f9b9e266567605e/formalization/distance-preservation/truncated12/TenGeometry.lean#L26-L87}{Complete checked declaration at the immutable proof commit.}

\lstinputlisting{listings/TenGeometry-source_ten_geometry.tex}

\subsection*{\texttt{TenEncoding}: ten\_geometry\_encodes}

\noindent\href{https://github.com/alectionary/vibe-math/blob/3ae21c0719cedbd928f4f7882f9b9e266567605e/formalization/distance-preservation/truncated12/TenEncoding.lean#L8-L38}{Complete checked declaration at the immutable proof commit.}

\lstinputlisting{listings/TenEncoding-ten_geometry_encodes.tex}

\subsection*{\texttt{TenEncoding}: actual\_canonical\_ten\_support\_encodes}

\noindent\href{https://github.com/alectionary/vibe-math/blob/3ae21c0719cedbd928f4f7882f9b9e266567605e/formalization/distance-preservation/truncated12/TenEncoding.lean#L40-L52}{Complete checked declaration at the immutable proof commit.}

\lstinputlisting{listings/TenEncoding-actual_canonical_ten_support_encodes.tex}

\subsection*{\texttt{K23Numbering}: actual\_k23\_numbering}

\noindent\href{https://github.com/alectionary/vibe-math/blob/3ae21c0719cedbd928f4f7882f9b9e266567605e/formalization/distance-preservation/truncated12/K23Numbering.lean#L29-L74}{Complete checked declaration at the immutable proof commit.}

\lstinputlisting{listings/K23Numbering-actual_k23_numbering.tex}

\subsection*{\texttt{ElevenGeometry}: actual\_eleven\_geometry}

\noindent\href{https://github.com/alectionary/vibe-math/blob/3ae21c0719cedbd928f4f7882f9b9e266567605e/formalization/distance-preservation/truncated12/ElevenGeometry.lean#L74-L115}{Complete checked declaration at the immutable proof commit.}

\lstinputlisting{listings/ElevenGeometry-actual_eleven_geometry.tex}

\subsection*{\texttt{SourceEleven}: source\_canonical\_eleven\_stabilizer}

\noindent\href{https://github.com/alectionary/vibe-math/blob/3ae21c0719cedbd928f4f7882f9b9e266567605e/formalization/distance-preservation/truncated12/SourceEleven.lean#L15-L51}{Complete checked declaration at the immutable proof commit.}

\lstinputlisting{listings/SourceEleven-source_canonical_eleven_stabilizer.tex}

\clearpage\section{Integer certificate interpretation and cyclic torsion}

Integer identities are interpreted in every additive group. The finite catalog supplies witnesses to this checked implication; it supplies no additional axiom.

\subsection*{\texttt{IntegerSoundness}: integer\_certificate\_sound}

\noindent\href{https://github.com/alectionary/vibe-math/blob/3ae21c0719cedbd928f4f7882f9b9e266567605e/formalization/distance-preservation/truncated12/IntegerSoundness.lean#L11-L22}{Complete checked declaration at the immutable proof commit.}

\lstinputlisting{listings/IntegerSoundness-integer_certificate_sound.tex}

\subsection*{\texttt{KernelCatalog}: checked\_fast\_equations}

\noindent\href{https://github.com/alectionary/vibe-math/blob/3ae21c0719cedbd928f4f7882f9b9e266567605e/formalization/distance-preservation/truncated12/KernelCatalog.lean#L41-L57}{Complete checked declaration at the immutable proof commit.}

\lstinputlisting{listings/KernelCatalog-checked_fast_equations.tex}

\subsection*{\texttt{PatternRealizations}: PatternRealization}

\noindent\href{https://github.com/alectionary/vibe-math/blob/3ae21c0719cedbd928f4f7882f9b9e266567605e/formalization/distance-preservation/truncated12/PatternRealizations.lean#L10-L21}{Complete checked declaration at the immutable proof commit.}

\lstinputlisting{listings/PatternRealizations-PatternRealization.tex}

\subsection*{\texttt{PatternRealizations}: certified\_pattern\_has\_no\_source\_realization}

\noindent\href{https://github.com/alectionary/vibe-math/blob/3ae21c0719cedbd928f4f7882f9b9e266567605e/formalization/distance-preservation/truncated12/PatternRealizations.lean#L54-L72}{Complete checked declaration at the immutable proof commit.}

\lstinputlisting{listings/PatternRealizations-certified_pattern_has_no_source_realization.tex}

\subsection*{\texttt{TorsionClassical}: torsionWord}

\noindent\href{https://github.com/alectionary/vibe-math/blob/3ae21c0719cedbd928f4f7882f9b9e266567605e/formalization/distance-preservation/truncated12/TorsionClassical.lean#L17-L19}{Complete checked declaration at the immutable proof commit.}

\lstinputlisting{listings/TorsionClassical-torsionWord.tex}

\subsection*{\texttt{TorsionClassical}: source\_classical\_word\_le10}

\noindent\href{https://github.com/alectionary/vibe-math/blob/3ae21c0719cedbd928f4f7882f9b9e266567605e/formalization/distance-preservation/truncated12/TorsionClassical.lean#L66-L74}{Complete checked declaration at the immutable proof commit.}

\lstinputlisting{listings/TorsionClassical-source_classical_word_le10.tex}

\subsection*{\texttt{CoverageCore}: Covered}

\noindent\href{https://github.com/alectionary/vibe-math/blob/3ae21c0719cedbd928f4f7882f9b9e266567605e/formalization/distance-preservation/truncated12/CoverageCore.lean#L7-L8}{Complete checked declaration at the immutable proof commit.}

\lstinputlisting{listings/CoverageCore-Covered.tex}

\subsection*{\texttt{CoverageCore}: covered\_single}

\noindent\href{https://github.com/alectionary/vibe-math/blob/3ae21c0719cedbd928f4f7882f9b9e266567605e/formalization/distance-preservation/truncated12/CoverageCore.lean#L10-L15}{Complete checked declaration at the immutable proof commit.}

\lstinputlisting{listings/CoverageCore-covered_single.tex}

\subsection*{\texttt{CoverageCore}: covered\_append}

\noindent\href{https://github.com/alectionary/vibe-math/blob/3ae21c0719cedbd928f4f7882f9b9e266567605e/formalization/distance-preservation/truncated12/CoverageCore.lean#L17-L26}{Complete checked declaration at the immutable proof commit.}

\lstinputlisting{listings/CoverageCore-covered_append.tex}

\clearpage\section{Transport and the complete global reduction}

Canonical group placement is discharged by permutation and re-normalization. The global proof below includes all three profiles and the pure sectors.

\subsection*{\texttt{ProfileTransport}: DistanceCriterion}

\noindent\href{https://github.com/alectionary/vibe-math/blob/3ae21c0719cedbd928f4f7882f9b9e266567605e/formalization/distance-preservation/truncated12/ProfileTransport.lean#L9-L17}{Complete checked declaration at the immutable proof commit.}

\lstinputlisting{listings/ProfileTransport-DistanceCriterion.tex}

\subsection*{\texttt{ProfileTransport}: normPerm\_criterion}

\noindent\href{https://github.com/alectionary/vibe-math/blob/3ae21c0719cedbd928f4f7882f9b9e266567605e/formalization/distance-preservation/truncated12/ProfileTransport.lean#L30-L43}{Complete checked declaration at the immutable proof commit.}

\lstinputlisting{listings/ProfileTransport-normPerm_criterion.tex}

\subsection*{\texttt{ProfileTransport}: normPerm\_span}

\noindent\href{https://github.com/alectionary/vibe-math/blob/3ae21c0719cedbd928f4f7882f9b9e266567605e/formalization/distance-preservation/truncated12/ProfileTransport.lean#L56-L64}{Complete checked declaration at the immutable proof commit.}

\lstinputlisting{listings/ProfileTransport-normPerm_span.tex}

\subsection*{\texttt{TwelveConditional}: source\_kernel\_below\_twelve\_of\_coverage}

\noindent\href{https://github.com/alectionary/vibe-math/blob/3ae21c0719cedbd928f4f7882f9b9e266567605e/formalization/distance-preservation/truncated12/TwelveConditional.lean#L13-L133}{Complete checked declaration at the immutable proof commit.}

\lstinputlisting{listings/TwelveConditional-source_kernel_below_twelve_of_coverage.tex}

\clearpage\section{Pauli union weight and the literal minimum theorem}

Both sectors and the entire product of stabilizer spans are retained. The public theorem discharges finite coverage before passing to literal minima.

\subsection*{\texttt{PauliBridge}: pauli\_bound\_of\_x\_bound}

\noindent\href{https://github.com/alectionary/vibe-math/blob/3ae21c0719cedbd928f4f7882f9b9e266567605e/formalization/distance-preservation/truncated12/PauliBridge.lean#L31-L42}{Complete checked declaration at the immutable proof commit.}

\lstinputlisting{listings/PauliBridge-pauli_bound_of_x_bound.tex}

\subsection*{\texttt{ThroughEleven}: source\_pauli\_bound\_twelve}

\noindent\href{https://github.com/alectionary/vibe-math/blob/3ae21c0719cedbd928f4f7882f9b9e266567605e/formalization/distance-preservation/truncated12/ThroughEleven.lean#L27-L39}{Complete checked declaration at the immutable proof commit.}

\lstinputlisting{listings/ThroughEleven-source_pauli_bound_twelve.tex}

\subsection*{\texttt{DistanceMinima}: source\_classical\_minimum\_realized}

\noindent\href{https://github.com/alectionary/vibe-math/blob/3ae21c0719cedbd928f4f7882f9b9e266567605e/formalization/distance-preservation/truncated12/DistanceMinima.lean#L24-L26}{Complete checked declaration at the immutable proof commit.}

\lstinputlisting{listings/DistanceMinima-source_classical_minimum_realized.tex}

\subsection*{\texttt{DistanceMinima}: source\_classical\_minimum\_lower}

\noindent\href{https://github.com/alectionary/vibe-math/blob/3ae21c0719cedbd928f4f7882f9b9e266567605e/formalization/distance-preservation/truncated12/DistanceMinima.lean#L28-L31}{Complete checked declaration at the immutable proof commit.}

\lstinputlisting{listings/DistanceMinima-source_classical_minimum_lower.tex}

\subsection*{\texttt{DistanceMinima}: source\_transpose\_zero\_kernel\_infinity}

\noindent\href{https://github.com/alectionary/vibe-math/blob/3ae21c0719cedbd928f4f7882f9b9e266567605e/formalization/distance-preservation/truncated12/DistanceMinima.lean#L40-L51}{Complete checked declaration at the immutable proof commit.}

\lstinputlisting{listings/DistanceMinima-source_transpose_zero_kernel_infinity.tex}

\subsection*{\texttt{DistanceMinima}: source\_quantum\_minimum\_lower\_iff}

\noindent\href{https://github.com/alectionary/vibe-math/blob/3ae21c0719cedbd928f4f7882f9b9e266567605e/formalization/distance-preservation/truncated12/DistanceMinima.lean#L88-L103}{Complete checked declaration at the immutable proof commit.}

\lstinputlisting{listings/DistanceMinima-source_quantum_minimum_lower_iff.tex}

\subsection*{\texttt{CriterionScope}: no\_full\_row\_of\_no\_triple}

\noindent\href{https://github.com/alectionary/vibe-math/blob/3ae21c0719cedbd928f4f7882f9b9e266567605e/formalization/distance-preservation/truncated12/CriterionScope.lean#L12-L17}{Complete checked declaration at the immutable proof commit.}

\lstinputlisting{listings/CriterionScope-no_full_row_of_no_triple.tex}

\subsection*{\texttt{DistanceTwelve}: source\_distance\_bound\_twelve}

\noindent\href{https://github.com/alectionary/vibe-math/blob/3ae21c0719cedbd928f4f7882f9b9e266567605e/formalization/distance-preservation/truncated12/DistanceTwelve.lean#L9-L18}{Complete checked declaration at the immutable proof commit.}

\lstinputlisting{listings/DistanceTwelve-source_distance_bound_twelve.tex}

\clearpage\section{Complete research-source link index}\label{app:links}
Every module in the final theorem and source-scope audit is listed below. Each row has a link to the master branch and an immutable link at commit \nolinkurl{3ae21c0719cedbd928f4f7882f9b9e266567605e}. The source archive includes these exact files and their matching build records. The accompanying HTML guide provides hashes and imports.
\footnotesize\begin{longtable}{@{}P{.62\textwidth}P{.14\textwidth}P{.14\textwidth}@{}}
\toprule Module & Current & Immutable\\\midrule\endhead
truncated12/PackedTables.lean & \href{https://github.com/alectionary/vibe-math/blob/master/formalization/distance-preservation/truncated12/PackedTables.lean}{source} & \href{https://github.com/alectionary/vibe-math/blob/3ae21c0719cedbd928f4f7882f9b9e266567605e/formalization/distance-preservation/truncated12/PackedTables.lean}{source}\\
truncated12/CatalogChecker.lean & \href{https://github.com/alectionary/vibe-math/blob/master/formalization/distance-preservation/truncated12/CatalogChecker.lean}{source} & \href{https://github.com/alectionary/vibe-math/blob/3ae21c0719cedbd928f4f7882f9b9e266567605e/formalization/distance-preservation/truncated12/CatalogChecker.lean}{source}\\
truncated12/SparseCatalogChecker.lean & \href{https://github.com/alectionary/vibe-math/blob/master/formalization/distance-preservation/truncated12/SparseCatalogChecker.lean}{source} & \href{https://github.com/alectionary/vibe-math/blob/3ae21c0719cedbd928f4f7882f9b9e266567605e/formalization/distance-preservation/truncated12/SparseCatalogChecker.lean}{source}\\
truncated12/FastCatalogChecker.lean & \href{https://github.com/alectionary/vibe-math/blob/master/formalization/distance-preservation/truncated12/FastCatalogChecker.lean}{source} & \href{https://github.com/alectionary/vibe-math/blob/3ae21c0719cedbd928f4f7882f9b9e266567605e/formalization/distance-preservation/truncated12/FastCatalogChecker.lean}{source}\\
truncated12/catalog-chunks/F10\_000.lean & \href{https://github.com/alectionary/vibe-math/blob/master/formalization/distance-preservation/truncated12/catalog-chunks/F10_000.lean}{source} & \href{https://github.com/alectionary/vibe-math/blob/3ae21c0719cedbd928f4f7882f9b9e266567605e/formalization/distance-preservation/truncated12/catalog-chunks/F10_000.lean}{source}\\
truncated12/catalog-chunks/F10\_001.lean & \href{https://github.com/alectionary/vibe-math/blob/master/formalization/distance-preservation/truncated12/catalog-chunks/F10_001.lean}{source} & \href{https://github.com/alectionary/vibe-math/blob/3ae21c0719cedbd928f4f7882f9b9e266567605e/formalization/distance-preservation/truncated12/catalog-chunks/F10_001.lean}{source}\\
truncated12/catalog-chunks/F10\_002.lean & \href{https://github.com/alectionary/vibe-math/blob/master/formalization/distance-preservation/truncated12/catalog-chunks/F10_002.lean}{source} & \href{https://github.com/alectionary/vibe-math/blob/3ae21c0719cedbd928f4f7882f9b9e266567605e/formalization/distance-preservation/truncated12/catalog-chunks/F10_002.lean}{source}\\
truncated12/catalog-chunks/F10\_003.lean & \href{https://github.com/alectionary/vibe-math/blob/master/formalization/distance-preservation/truncated12/catalog-chunks/F10_003.lean}{source} & \href{https://github.com/alectionary/vibe-math/blob/3ae21c0719cedbd928f4f7882f9b9e266567605e/formalization/distance-preservation/truncated12/catalog-chunks/F10_003.lean}{source}\\
truncated12/catalog-chunks/F10\_004.lean & \href{https://github.com/alectionary/vibe-math/blob/master/formalization/distance-preservation/truncated12/catalog-chunks/F10_004.lean}{source} & \href{https://github.com/alectionary/vibe-math/blob/3ae21c0719cedbd928f4f7882f9b9e266567605e/formalization/distance-preservation/truncated12/catalog-chunks/F10_004.lean}{source}\\
truncated12/catalog-chunks/F10\_005.lean & \href{https://github.com/alectionary/vibe-math/blob/master/formalization/distance-preservation/truncated12/catalog-chunks/F10_005.lean}{source} & \href{https://github.com/alectionary/vibe-math/blob/3ae21c0719cedbd928f4f7882f9b9e266567605e/formalization/distance-preservation/truncated12/catalog-chunks/F10_005.lean}{source}\\
truncated12/catalog-chunks/F10\_006.lean & \href{https://github.com/alectionary/vibe-math/blob/master/formalization/distance-preservation/truncated12/catalog-chunks/F10_006.lean}{source} & \href{https://github.com/alectionary/vibe-math/blob/3ae21c0719cedbd928f4f7882f9b9e266567605e/formalization/distance-preservation/truncated12/catalog-chunks/F10_006.lean}{source}\\
truncated12/catalog-chunks/F10\_007.lean & \href{https://github.com/alectionary/vibe-math/blob/master/formalization/distance-preservation/truncated12/catalog-chunks/F10_007.lean}{source} & \href{https://github.com/alectionary/vibe-math/blob/3ae21c0719cedbd928f4f7882f9b9e266567605e/formalization/distance-preservation/truncated12/catalog-chunks/F10_007.lean}{source}\\
truncated12/catalog-chunks/F10\_008.lean & \href{https://github.com/alectionary/vibe-math/blob/master/formalization/distance-preservation/truncated12/catalog-chunks/F10_008.lean}{source} & \href{https://github.com/alectionary/vibe-math/blob/3ae21c0719cedbd928f4f7882f9b9e266567605e/formalization/distance-preservation/truncated12/catalog-chunks/F10_008.lean}{source}\\
truncated12/catalog-chunks/F10\_009.lean & \href{https://github.com/alectionary/vibe-math/blob/master/formalization/distance-preservation/truncated12/catalog-chunks/F10_009.lean}{source} & \href{https://github.com/alectionary/vibe-math/blob/3ae21c0719cedbd928f4f7882f9b9e266567605e/formalization/distance-preservation/truncated12/catalog-chunks/F10_009.lean}{source}\\
truncated12/catalog-chunks/F10\_010.lean & \href{https://github.com/alectionary/vibe-math/blob/master/formalization/distance-preservation/truncated12/catalog-chunks/F10_010.lean}{source} & \href{https://github.com/alectionary/vibe-math/blob/3ae21c0719cedbd928f4f7882f9b9e266567605e/formalization/distance-preservation/truncated12/catalog-chunks/F10_010.lean}{source}\\
truncated12/catalog-chunks/F10\_011.lean & \href{https://github.com/alectionary/vibe-math/blob/master/formalization/distance-preservation/truncated12/catalog-chunks/F10_011.lean}{source} & \href{https://github.com/alectionary/vibe-math/blob/3ae21c0719cedbd928f4f7882f9b9e266567605e/formalization/distance-preservation/truncated12/catalog-chunks/F10_011.lean}{source}\\
truncated12/catalog-chunks/F10\_012.lean & \href{https://github.com/alectionary/vibe-math/blob/master/formalization/distance-preservation/truncated12/catalog-chunks/F10_012.lean}{source} & \href{https://github.com/alectionary/vibe-math/blob/3ae21c0719cedbd928f4f7882f9b9e266567605e/formalization/distance-preservation/truncated12/catalog-chunks/F10_012.lean}{source}\\
truncated12/catalog-chunks/F10\_013.lean & \href{https://github.com/alectionary/vibe-math/blob/master/formalization/distance-preservation/truncated12/catalog-chunks/F10_013.lean}{source} & \href{https://github.com/alectionary/vibe-math/blob/3ae21c0719cedbd928f4f7882f9b9e266567605e/formalization/distance-preservation/truncated12/catalog-chunks/F10_013.lean}{source}\\
truncated12/catalog-chunks/F10\_014.lean & \href{https://github.com/alectionary/vibe-math/blob/master/formalization/distance-preservation/truncated12/catalog-chunks/F10_014.lean}{source} & \href{https://github.com/alectionary/vibe-math/blob/3ae21c0719cedbd928f4f7882f9b9e266567605e/formalization/distance-preservation/truncated12/catalog-chunks/F10_014.lean}{source}\\
truncated12/catalog-chunks/F10\_015.lean & \href{https://github.com/alectionary/vibe-math/blob/master/formalization/distance-preservation/truncated12/catalog-chunks/F10_015.lean}{source} & \href{https://github.com/alectionary/vibe-math/blob/3ae21c0719cedbd928f4f7882f9b9e266567605e/formalization/distance-preservation/truncated12/catalog-chunks/F10_015.lean}{source}\\
truncated12/catalog-chunks/F10\_016.lean & \href{https://github.com/alectionary/vibe-math/blob/master/formalization/distance-preservation/truncated12/catalog-chunks/F10_016.lean}{source} & \href{https://github.com/alectionary/vibe-math/blob/3ae21c0719cedbd928f4f7882f9b9e266567605e/formalization/distance-preservation/truncated12/catalog-chunks/F10_016.lean}{source}\\
truncated12/catalog-chunks/F10\_017.lean & \href{https://github.com/alectionary/vibe-math/blob/master/formalization/distance-preservation/truncated12/catalog-chunks/F10_017.lean}{source} & \href{https://github.com/alectionary/vibe-math/blob/3ae21c0719cedbd928f4f7882f9b9e266567605e/formalization/distance-preservation/truncated12/catalog-chunks/F10_017.lean}{source}\\
truncated12/catalog-chunks/F10\_018.lean & \href{https://github.com/alectionary/vibe-math/blob/master/formalization/distance-preservation/truncated12/catalog-chunks/F10_018.lean}{source} & \href{https://github.com/alectionary/vibe-math/blob/3ae21c0719cedbd928f4f7882f9b9e266567605e/formalization/distance-preservation/truncated12/catalog-chunks/F10_018.lean}{source}\\
truncated12/catalog-chunks/F10\_019.lean & \href{https://github.com/alectionary/vibe-math/blob/master/formalization/distance-preservation/truncated12/catalog-chunks/F10_019.lean}{source} & \href{https://github.com/alectionary/vibe-math/blob/3ae21c0719cedbd928f4f7882f9b9e266567605e/formalization/distance-preservation/truncated12/catalog-chunks/F10_019.lean}{source}\\
truncated12/catalog-chunks/F10\_020.lean & \href{https://github.com/alectionary/vibe-math/blob/master/formalization/distance-preservation/truncated12/catalog-chunks/F10_020.lean}{source} & \href{https://github.com/alectionary/vibe-math/blob/3ae21c0719cedbd928f4f7882f9b9e266567605e/formalization/distance-preservation/truncated12/catalog-chunks/F10_020.lean}{source}\\
truncated12/catalog-chunks/F10\_021.lean & \href{https://github.com/alectionary/vibe-math/blob/master/formalization/distance-preservation/truncated12/catalog-chunks/F10_021.lean}{source} & \href{https://github.com/alectionary/vibe-math/blob/3ae21c0719cedbd928f4f7882f9b9e266567605e/formalization/distance-preservation/truncated12/catalog-chunks/F10_021.lean}{source}\\
truncated12/catalog-chunks/F10\_022.lean & \href{https://github.com/alectionary/vibe-math/blob/master/formalization/distance-preservation/truncated12/catalog-chunks/F10_022.lean}{source} & \href{https://github.com/alectionary/vibe-math/blob/3ae21c0719cedbd928f4f7882f9b9e266567605e/formalization/distance-preservation/truncated12/catalog-chunks/F10_022.lean}{source}\\
truncated12/catalog-chunks/F10\_023.lean & \href{https://github.com/alectionary/vibe-math/blob/master/formalization/distance-preservation/truncated12/catalog-chunks/F10_023.lean}{source} & \href{https://github.com/alectionary/vibe-math/blob/3ae21c0719cedbd928f4f7882f9b9e266567605e/formalization/distance-preservation/truncated12/catalog-chunks/F10_023.lean}{source}\\
truncated12/catalog-chunks/F10\_024.lean & \href{https://github.com/alectionary/vibe-math/blob/master/formalization/distance-preservation/truncated12/catalog-chunks/F10_024.lean}{source} & \href{https://github.com/alectionary/vibe-math/blob/3ae21c0719cedbd928f4f7882f9b9e266567605e/formalization/distance-preservation/truncated12/catalog-chunks/F10_024.lean}{source}\\
truncated12/catalog-chunks/F10\_025.lean & \href{https://github.com/alectionary/vibe-math/blob/master/formalization/distance-preservation/truncated12/catalog-chunks/F10_025.lean}{source} & \href{https://github.com/alectionary/vibe-math/blob/3ae21c0719cedbd928f4f7882f9b9e266567605e/formalization/distance-preservation/truncated12/catalog-chunks/F10_025.lean}{source}\\
truncated12/catalog-chunks/F10\_026.lean & \href{https://github.com/alectionary/vibe-math/blob/master/formalization/distance-preservation/truncated12/catalog-chunks/F10_026.lean}{source} & \href{https://github.com/alectionary/vibe-math/blob/3ae21c0719cedbd928f4f7882f9b9e266567605e/formalization/distance-preservation/truncated12/catalog-chunks/F10_026.lean}{source}\\
truncated12/catalog-chunks/F10\_027.lean & \href{https://github.com/alectionary/vibe-math/blob/master/formalization/distance-preservation/truncated12/catalog-chunks/F10_027.lean}{source} & \href{https://github.com/alectionary/vibe-math/blob/3ae21c0719cedbd928f4f7882f9b9e266567605e/formalization/distance-preservation/truncated12/catalog-chunks/F10_027.lean}{source}\\
truncated12/catalog-chunks/F10\_028.lean & \href{https://github.com/alectionary/vibe-math/blob/master/formalization/distance-preservation/truncated12/catalog-chunks/F10_028.lean}{source} & \href{https://github.com/alectionary/vibe-math/blob/3ae21c0719cedbd928f4f7882f9b9e266567605e/formalization/distance-preservation/truncated12/catalog-chunks/F10_028.lean}{source}\\
truncated12/catalog-chunks/F10\_029.lean & \href{https://github.com/alectionary/vibe-math/blob/master/formalization/distance-preservation/truncated12/catalog-chunks/F10_029.lean}{source} & \href{https://github.com/alectionary/vibe-math/blob/3ae21c0719cedbd928f4f7882f9b9e266567605e/formalization/distance-preservation/truncated12/catalog-chunks/F10_029.lean}{source}\\
truncated12/catalog-chunks/F10\_030.lean & \href{https://github.com/alectionary/vibe-math/blob/master/formalization/distance-preservation/truncated12/catalog-chunks/F10_030.lean}{source} & \href{https://github.com/alectionary/vibe-math/blob/3ae21c0719cedbd928f4f7882f9b9e266567605e/formalization/distance-preservation/truncated12/catalog-chunks/F10_030.lean}{source}\\
truncated12/catalog-chunks/F10\_031.lean & \href{https://github.com/alectionary/vibe-math/blob/master/formalization/distance-preservation/truncated12/catalog-chunks/F10_031.lean}{source} & \href{https://github.com/alectionary/vibe-math/blob/3ae21c0719cedbd928f4f7882f9b9e266567605e/formalization/distance-preservation/truncated12/catalog-chunks/F10_031.lean}{source}\\
truncated12/catalog-chunks/F10\_032.lean & \href{https://github.com/alectionary/vibe-math/blob/master/formalization/distance-preservation/truncated12/catalog-chunks/F10_032.lean}{source} & \href{https://github.com/alectionary/vibe-math/blob/3ae21c0719cedbd928f4f7882f9b9e266567605e/formalization/distance-preservation/truncated12/catalog-chunks/F10_032.lean}{source}\\
truncated12/catalog-chunks/F10\_033.lean & \href{https://github.com/alectionary/vibe-math/blob/master/formalization/distance-preservation/truncated12/catalog-chunks/F10_033.lean}{source} & \href{https://github.com/alectionary/vibe-math/blob/3ae21c0719cedbd928f4f7882f9b9e266567605e/formalization/distance-preservation/truncated12/catalog-chunks/F10_033.lean}{source}\\
truncated12/catalog-chunks/F10\_034.lean & \href{https://github.com/alectionary/vibe-math/blob/master/formalization/distance-preservation/truncated12/catalog-chunks/F10_034.lean}{source} & \href{https://github.com/alectionary/vibe-math/blob/3ae21c0719cedbd928f4f7882f9b9e266567605e/formalization/distance-preservation/truncated12/catalog-chunks/F10_034.lean}{source}\\
truncated12/catalog-chunks/F10\_035.lean & \href{https://github.com/alectionary/vibe-math/blob/master/formalization/distance-preservation/truncated12/catalog-chunks/F10_035.lean}{source} & \href{https://github.com/alectionary/vibe-math/blob/3ae21c0719cedbd928f4f7882f9b9e266567605e/formalization/distance-preservation/truncated12/catalog-chunks/F10_035.lean}{source}\\
truncated12/catalog-chunks/F10\_036.lean & \href{https://github.com/alectionary/vibe-math/blob/master/formalization/distance-preservation/truncated12/catalog-chunks/F10_036.lean}{source} & \href{https://github.com/alectionary/vibe-math/blob/3ae21c0719cedbd928f4f7882f9b9e266567605e/formalization/distance-preservation/truncated12/catalog-chunks/F10_036.lean}{source}\\
truncated12/catalog-chunks/F10\_037.lean & \href{https://github.com/alectionary/vibe-math/blob/master/formalization/distance-preservation/truncated12/catalog-chunks/F10_037.lean}{source} & \href{https://github.com/alectionary/vibe-math/blob/3ae21c0719cedbd928f4f7882f9b9e266567605e/formalization/distance-preservation/truncated12/catalog-chunks/F10_037.lean}{source}\\
truncated12/catalog-chunks/F10\_038.lean & \href{https://github.com/alectionary/vibe-math/blob/master/formalization/distance-preservation/truncated12/catalog-chunks/F10_038.lean}{source} & \href{https://github.com/alectionary/vibe-math/blob/3ae21c0719cedbd928f4f7882f9b9e266567605e/formalization/distance-preservation/truncated12/catalog-chunks/F10_038.lean}{source}\\
truncated12/catalog-chunks/F10\_039.lean & \href{https://github.com/alectionary/vibe-math/blob/master/formalization/distance-preservation/truncated12/catalog-chunks/F10_039.lean}{source} & \href{https://github.com/alectionary/vibe-math/blob/3ae21c0719cedbd928f4f7882f9b9e266567605e/formalization/distance-preservation/truncated12/catalog-chunks/F10_039.lean}{source}\\
truncated12/catalog-chunks/F10\_040.lean & \href{https://github.com/alectionary/vibe-math/blob/master/formalization/distance-preservation/truncated12/catalog-chunks/F10_040.lean}{source} & \href{https://github.com/alectionary/vibe-math/blob/3ae21c0719cedbd928f4f7882f9b9e266567605e/formalization/distance-preservation/truncated12/catalog-chunks/F10_040.lean}{source}\\
truncated12/catalog-chunks/F10\_041.lean & \href{https://github.com/alectionary/vibe-math/blob/master/formalization/distance-preservation/truncated12/catalog-chunks/F10_041.lean}{source} & \href{https://github.com/alectionary/vibe-math/blob/3ae21c0719cedbd928f4f7882f9b9e266567605e/formalization/distance-preservation/truncated12/catalog-chunks/F10_041.lean}{source}\\
truncated12/catalog-chunks/F10\_042.lean & \href{https://github.com/alectionary/vibe-math/blob/master/formalization/distance-preservation/truncated12/catalog-chunks/F10_042.lean}{source} & \href{https://github.com/alectionary/vibe-math/blob/3ae21c0719cedbd928f4f7882f9b9e266567605e/formalization/distance-preservation/truncated12/catalog-chunks/F10_042.lean}{source}\\
truncated12/catalog-chunks/F10\_043.lean & \href{https://github.com/alectionary/vibe-math/blob/master/formalization/distance-preservation/truncated12/catalog-chunks/F10_043.lean}{source} & \href{https://github.com/alectionary/vibe-math/blob/3ae21c0719cedbd928f4f7882f9b9e266567605e/formalization/distance-preservation/truncated12/catalog-chunks/F10_043.lean}{source}\\
truncated12/catalog-chunks/F10\_044.lean & \href{https://github.com/alectionary/vibe-math/blob/master/formalization/distance-preservation/truncated12/catalog-chunks/F10_044.lean}{source} & \href{https://github.com/alectionary/vibe-math/blob/3ae21c0719cedbd928f4f7882f9b9e266567605e/formalization/distance-preservation/truncated12/catalog-chunks/F10_044.lean}{source}\\
truncated12/catalog-chunks/F10\_045.lean & \href{https://github.com/alectionary/vibe-math/blob/master/formalization/distance-preservation/truncated12/catalog-chunks/F10_045.lean}{source} & \href{https://github.com/alectionary/vibe-math/blob/3ae21c0719cedbd928f4f7882f9b9e266567605e/formalization/distance-preservation/truncated12/catalog-chunks/F10_045.lean}{source}\\
truncated12/catalog-chunks/F10\_046.lean & \href{https://github.com/alectionary/vibe-math/blob/master/formalization/distance-preservation/truncated12/catalog-chunks/F10_046.lean}{source} & \href{https://github.com/alectionary/vibe-math/blob/3ae21c0719cedbd928f4f7882f9b9e266567605e/formalization/distance-preservation/truncated12/catalog-chunks/F10_046.lean}{source}\\
truncated12/catalog-chunks/F10\_047.lean & \href{https://github.com/alectionary/vibe-math/blob/master/formalization/distance-preservation/truncated12/catalog-chunks/F10_047.lean}{source} & \href{https://github.com/alectionary/vibe-math/blob/3ae21c0719cedbd928f4f7882f9b9e266567605e/formalization/distance-preservation/truncated12/catalog-chunks/F10_047.lean}{source}\\
truncated12/catalog-chunks/F10\_048.lean & \href{https://github.com/alectionary/vibe-math/blob/master/formalization/distance-preservation/truncated12/catalog-chunks/F10_048.lean}{source} & \href{https://github.com/alectionary/vibe-math/blob/3ae21c0719cedbd928f4f7882f9b9e266567605e/formalization/distance-preservation/truncated12/catalog-chunks/F10_048.lean}{source}\\
truncated12/catalog-chunks/F10\_049.lean & \href{https://github.com/alectionary/vibe-math/blob/master/formalization/distance-preservation/truncated12/catalog-chunks/F10_049.lean}{source} & \href{https://github.com/alectionary/vibe-math/blob/3ae21c0719cedbd928f4f7882f9b9e266567605e/formalization/distance-preservation/truncated12/catalog-chunks/F10_049.lean}{source}\\
truncated12/catalog-chunks/F10\_050.lean & \href{https://github.com/alectionary/vibe-math/blob/master/formalization/distance-preservation/truncated12/catalog-chunks/F10_050.lean}{source} & \href{https://github.com/alectionary/vibe-math/blob/3ae21c0719cedbd928f4f7882f9b9e266567605e/formalization/distance-preservation/truncated12/catalog-chunks/F10_050.lean}{source}\\
truncated12/catalog-chunks/F10\_051.lean & \href{https://github.com/alectionary/vibe-math/blob/master/formalization/distance-preservation/truncated12/catalog-chunks/F10_051.lean}{source} & \href{https://github.com/alectionary/vibe-math/blob/3ae21c0719cedbd928f4f7882f9b9e266567605e/formalization/distance-preservation/truncated12/catalog-chunks/F10_051.lean}{source}\\
truncated12/catalog-chunks/F10\_052.lean & \href{https://github.com/alectionary/vibe-math/blob/master/formalization/distance-preservation/truncated12/catalog-chunks/F10_052.lean}{source} & \href{https://github.com/alectionary/vibe-math/blob/3ae21c0719cedbd928f4f7882f9b9e266567605e/formalization/distance-preservation/truncated12/catalog-chunks/F10_052.lean}{source}\\
truncated12/catalog-chunks/F10\_053.lean & \href{https://github.com/alectionary/vibe-math/blob/master/formalization/distance-preservation/truncated12/catalog-chunks/F10_053.lean}{source} & \href{https://github.com/alectionary/vibe-math/blob/3ae21c0719cedbd928f4f7882f9b9e266567605e/formalization/distance-preservation/truncated12/catalog-chunks/F10_053.lean}{source}\\
truncated12/catalog-chunks/F10\_054.lean & \href{https://github.com/alectionary/vibe-math/blob/master/formalization/distance-preservation/truncated12/catalog-chunks/F10_054.lean}{source} & \href{https://github.com/alectionary/vibe-math/blob/3ae21c0719cedbd928f4f7882f9b9e266567605e/formalization/distance-preservation/truncated12/catalog-chunks/F10_054.lean}{source}\\
truncated12/catalog-chunks/F10\_055.lean & \href{https://github.com/alectionary/vibe-math/blob/master/formalization/distance-preservation/truncated12/catalog-chunks/F10_055.lean}{source} & \href{https://github.com/alectionary/vibe-math/blob/3ae21c0719cedbd928f4f7882f9b9e266567605e/formalization/distance-preservation/truncated12/catalog-chunks/F10_055.lean}{source}\\
truncated12/catalog-chunks/F10\_056.lean & \href{https://github.com/alectionary/vibe-math/blob/master/formalization/distance-preservation/truncated12/catalog-chunks/F10_056.lean}{source} & \href{https://github.com/alectionary/vibe-math/blob/3ae21c0719cedbd928f4f7882f9b9e266567605e/formalization/distance-preservation/truncated12/catalog-chunks/F10_056.lean}{source}\\
truncated12/catalog-chunks/F10\_057.lean & \href{https://github.com/alectionary/vibe-math/blob/master/formalization/distance-preservation/truncated12/catalog-chunks/F10_057.lean}{source} & \href{https://github.com/alectionary/vibe-math/blob/3ae21c0719cedbd928f4f7882f9b9e266567605e/formalization/distance-preservation/truncated12/catalog-chunks/F10_057.lean}{source}\\
truncated12/catalog-chunks/F10\_058.lean & \href{https://github.com/alectionary/vibe-math/blob/master/formalization/distance-preservation/truncated12/catalog-chunks/F10_058.lean}{source} & \href{https://github.com/alectionary/vibe-math/blob/3ae21c0719cedbd928f4f7882f9b9e266567605e/formalization/distance-preservation/truncated12/catalog-chunks/F10_058.lean}{source}\\
truncated12/catalog-chunks/F10\_059.lean & \href{https://github.com/alectionary/vibe-math/blob/master/formalization/distance-preservation/truncated12/catalog-chunks/F10_059.lean}{source} & \href{https://github.com/alectionary/vibe-math/blob/3ae21c0719cedbd928f4f7882f9b9e266567605e/formalization/distance-preservation/truncated12/catalog-chunks/F10_059.lean}{source}\\
truncated12/catalog-chunks/F10\_060.lean & \href{https://github.com/alectionary/vibe-math/blob/master/formalization/distance-preservation/truncated12/catalog-chunks/F10_060.lean}{source} & \href{https://github.com/alectionary/vibe-math/blob/3ae21c0719cedbd928f4f7882f9b9e266567605e/formalization/distance-preservation/truncated12/catalog-chunks/F10_060.lean}{source}\\
truncated12/catalog-chunks/F10\_061.lean & \href{https://github.com/alectionary/vibe-math/blob/master/formalization/distance-preservation/truncated12/catalog-chunks/F10_061.lean}{source} & \href{https://github.com/alectionary/vibe-math/blob/3ae21c0719cedbd928f4f7882f9b9e266567605e/formalization/distance-preservation/truncated12/catalog-chunks/F10_061.lean}{source}\\
truncated12/catalog-chunks/F10\_062.lean & \href{https://github.com/alectionary/vibe-math/blob/master/formalization/distance-preservation/truncated12/catalog-chunks/F10_062.lean}{source} & \href{https://github.com/alectionary/vibe-math/blob/3ae21c0719cedbd928f4f7882f9b9e266567605e/formalization/distance-preservation/truncated12/catalog-chunks/F10_062.lean}{source}\\
truncated12/catalog-chunks/F10\_063.lean & \href{https://github.com/alectionary/vibe-math/blob/master/formalization/distance-preservation/truncated12/catalog-chunks/F10_063.lean}{source} & \href{https://github.com/alectionary/vibe-math/blob/3ae21c0719cedbd928f4f7882f9b9e266567605e/formalization/distance-preservation/truncated12/catalog-chunks/F10_063.lean}{source}\\
truncated12/catalog-chunks/F10\_064.lean & \href{https://github.com/alectionary/vibe-math/blob/master/formalization/distance-preservation/truncated12/catalog-chunks/F10_064.lean}{source} & \href{https://github.com/alectionary/vibe-math/blob/3ae21c0719cedbd928f4f7882f9b9e266567605e/formalization/distance-preservation/truncated12/catalog-chunks/F10_064.lean}{source}\\
truncated12/catalog-chunks/F10\_065.lean & \href{https://github.com/alectionary/vibe-math/blob/master/formalization/distance-preservation/truncated12/catalog-chunks/F10_065.lean}{source} & \href{https://github.com/alectionary/vibe-math/blob/3ae21c0719cedbd928f4f7882f9b9e266567605e/formalization/distance-preservation/truncated12/catalog-chunks/F10_065.lean}{source}\\
truncated12/catalog-chunks/F10\_066.lean & \href{https://github.com/alectionary/vibe-math/blob/master/formalization/distance-preservation/truncated12/catalog-chunks/F10_066.lean}{source} & \href{https://github.com/alectionary/vibe-math/blob/3ae21c0719cedbd928f4f7882f9b9e266567605e/formalization/distance-preservation/truncated12/catalog-chunks/F10_066.lean}{source}\\
truncated12/catalog-chunks/F10\_067.lean & \href{https://github.com/alectionary/vibe-math/blob/master/formalization/distance-preservation/truncated12/catalog-chunks/F10_067.lean}{source} & \href{https://github.com/alectionary/vibe-math/blob/3ae21c0719cedbd928f4f7882f9b9e266567605e/formalization/distance-preservation/truncated12/catalog-chunks/F10_067.lean}{source}\\
truncated12/catalog-chunks/F10\_068.lean & \href{https://github.com/alectionary/vibe-math/blob/master/formalization/distance-preservation/truncated12/catalog-chunks/F10_068.lean}{source} & \href{https://github.com/alectionary/vibe-math/blob/3ae21c0719cedbd928f4f7882f9b9e266567605e/formalization/distance-preservation/truncated12/catalog-chunks/F10_068.lean}{source}\\
truncated12/catalog-chunks/F10\_069.lean & \href{https://github.com/alectionary/vibe-math/blob/master/formalization/distance-preservation/truncated12/catalog-chunks/F10_069.lean}{source} & \href{https://github.com/alectionary/vibe-math/blob/3ae21c0719cedbd928f4f7882f9b9e266567605e/formalization/distance-preservation/truncated12/catalog-chunks/F10_069.lean}{source}\\
truncated12/catalog-chunks/F10\_070.lean & \href{https://github.com/alectionary/vibe-math/blob/master/formalization/distance-preservation/truncated12/catalog-chunks/F10_070.lean}{source} & \href{https://github.com/alectionary/vibe-math/blob/3ae21c0719cedbd928f4f7882f9b9e266567605e/formalization/distance-preservation/truncated12/catalog-chunks/F10_070.lean}{source}\\
truncated12/catalog-chunks/F10\_071.lean & \href{https://github.com/alectionary/vibe-math/blob/master/formalization/distance-preservation/truncated12/catalog-chunks/F10_071.lean}{source} & \href{https://github.com/alectionary/vibe-math/blob/3ae21c0719cedbd928f4f7882f9b9e266567605e/formalization/distance-preservation/truncated12/catalog-chunks/F10_071.lean}{source}\\
truncated12/catalog-chunks/F10\_072.lean & \href{https://github.com/alectionary/vibe-math/blob/master/formalization/distance-preservation/truncated12/catalog-chunks/F10_072.lean}{source} & \href{https://github.com/alectionary/vibe-math/blob/3ae21c0719cedbd928f4f7882f9b9e266567605e/formalization/distance-preservation/truncated12/catalog-chunks/F10_072.lean}{source}\\
truncated12/catalog-chunks/F10\_073.lean & \href{https://github.com/alectionary/vibe-math/blob/master/formalization/distance-preservation/truncated12/catalog-chunks/F10_073.lean}{source} & \href{https://github.com/alectionary/vibe-math/blob/3ae21c0719cedbd928f4f7882f9b9e266567605e/formalization/distance-preservation/truncated12/catalog-chunks/F10_073.lean}{source}\\
truncated12/catalog-chunks/F10\_074.lean & \href{https://github.com/alectionary/vibe-math/blob/master/formalization/distance-preservation/truncated12/catalog-chunks/F10_074.lean}{source} & \href{https://github.com/alectionary/vibe-math/blob/3ae21c0719cedbd928f4f7882f9b9e266567605e/formalization/distance-preservation/truncated12/catalog-chunks/F10_074.lean}{source}\\
truncated12/catalog-chunks/F10\_075.lean & \href{https://github.com/alectionary/vibe-math/blob/master/formalization/distance-preservation/truncated12/catalog-chunks/F10_075.lean}{source} & \href{https://github.com/alectionary/vibe-math/blob/3ae21c0719cedbd928f4f7882f9b9e266567605e/formalization/distance-preservation/truncated12/catalog-chunks/F10_075.lean}{source}\\
truncated12/catalog-chunks/F10\_076.lean & \href{https://github.com/alectionary/vibe-math/blob/master/formalization/distance-preservation/truncated12/catalog-chunks/F10_076.lean}{source} & \href{https://github.com/alectionary/vibe-math/blob/3ae21c0719cedbd928f4f7882f9b9e266567605e/formalization/distance-preservation/truncated12/catalog-chunks/F10_076.lean}{source}\\
truncated12/catalog-chunks/F10\_077.lean & \href{https://github.com/alectionary/vibe-math/blob/master/formalization/distance-preservation/truncated12/catalog-chunks/F10_077.lean}{source} & \href{https://github.com/alectionary/vibe-math/blob/3ae21c0719cedbd928f4f7882f9b9e266567605e/formalization/distance-preservation/truncated12/catalog-chunks/F10_077.lean}{source}\\
truncated12/catalog-chunks/F10\_078.lean & \href{https://github.com/alectionary/vibe-math/blob/master/formalization/distance-preservation/truncated12/catalog-chunks/F10_078.lean}{source} & \href{https://github.com/alectionary/vibe-math/blob/3ae21c0719cedbd928f4f7882f9b9e266567605e/formalization/distance-preservation/truncated12/catalog-chunks/F10_078.lean}{source}\\
truncated12/catalog-chunks/F10\_079.lean & \href{https://github.com/alectionary/vibe-math/blob/master/formalization/distance-preservation/truncated12/catalog-chunks/F10_079.lean}{source} & \href{https://github.com/alectionary/vibe-math/blob/3ae21c0719cedbd928f4f7882f9b9e266567605e/formalization/distance-preservation/truncated12/catalog-chunks/F10_079.lean}{source}\\
truncated12/catalog-chunks/F10\_080.lean & \href{https://github.com/alectionary/vibe-math/blob/master/formalization/distance-preservation/truncated12/catalog-chunks/F10_080.lean}{source} & \href{https://github.com/alectionary/vibe-math/blob/3ae21c0719cedbd928f4f7882f9b9e266567605e/formalization/distance-preservation/truncated12/catalog-chunks/F10_080.lean}{source}\\
truncated12/catalog-chunks/F10\_081.lean & \href{https://github.com/alectionary/vibe-math/blob/master/formalization/distance-preservation/truncated12/catalog-chunks/F10_081.lean}{source} & \href{https://github.com/alectionary/vibe-math/blob/3ae21c0719cedbd928f4f7882f9b9e266567605e/formalization/distance-preservation/truncated12/catalog-chunks/F10_081.lean}{source}\\
truncated12/catalog-chunks/F10\_082.lean & \href{https://github.com/alectionary/vibe-math/blob/master/formalization/distance-preservation/truncated12/catalog-chunks/F10_082.lean}{source} & \href{https://github.com/alectionary/vibe-math/blob/3ae21c0719cedbd928f4f7882f9b9e266567605e/formalization/distance-preservation/truncated12/catalog-chunks/F10_082.lean}{source}\\
truncated12/catalog-chunks/F10\_083.lean & \href{https://github.com/alectionary/vibe-math/blob/master/formalization/distance-preservation/truncated12/catalog-chunks/F10_083.lean}{source} & \href{https://github.com/alectionary/vibe-math/blob/3ae21c0719cedbd928f4f7882f9b9e266567605e/formalization/distance-preservation/truncated12/catalog-chunks/F10_083.lean}{source}\\
truncated12/catalog-chunks/F10\_084.lean & \href{https://github.com/alectionary/vibe-math/blob/master/formalization/distance-preservation/truncated12/catalog-chunks/F10_084.lean}{source} & \href{https://github.com/alectionary/vibe-math/blob/3ae21c0719cedbd928f4f7882f9b9e266567605e/formalization/distance-preservation/truncated12/catalog-chunks/F10_084.lean}{source}\\
truncated12/catalog-chunks/F10\_085.lean & \href{https://github.com/alectionary/vibe-math/blob/master/formalization/distance-preservation/truncated12/catalog-chunks/F10_085.lean}{source} & \href{https://github.com/alectionary/vibe-math/blob/3ae21c0719cedbd928f4f7882f9b9e266567605e/formalization/distance-preservation/truncated12/catalog-chunks/F10_085.lean}{source}\\
truncated12/catalog-chunks/F10\_086.lean & \href{https://github.com/alectionary/vibe-math/blob/master/formalization/distance-preservation/truncated12/catalog-chunks/F10_086.lean}{source} & \href{https://github.com/alectionary/vibe-math/blob/3ae21c0719cedbd928f4f7882f9b9e266567605e/formalization/distance-preservation/truncated12/catalog-chunks/F10_086.lean}{source}\\
truncated12/catalog-chunks/F10\_087.lean & \href{https://github.com/alectionary/vibe-math/blob/master/formalization/distance-preservation/truncated12/catalog-chunks/F10_087.lean}{source} & \href{https://github.com/alectionary/vibe-math/blob/3ae21c0719cedbd928f4f7882f9b9e266567605e/formalization/distance-preservation/truncated12/catalog-chunks/F10_087.lean}{source}\\
truncated12/catalog-chunks/F10\_088.lean & \href{https://github.com/alectionary/vibe-math/blob/master/formalization/distance-preservation/truncated12/catalog-chunks/F10_088.lean}{source} & \href{https://github.com/alectionary/vibe-math/blob/3ae21c0719cedbd928f4f7882f9b9e266567605e/formalization/distance-preservation/truncated12/catalog-chunks/F10_088.lean}{source}\\
truncated12/catalog-chunks/F10\_089.lean & \href{https://github.com/alectionary/vibe-math/blob/master/formalization/distance-preservation/truncated12/catalog-chunks/F10_089.lean}{source} & \href{https://github.com/alectionary/vibe-math/blob/3ae21c0719cedbd928f4f7882f9b9e266567605e/formalization/distance-preservation/truncated12/catalog-chunks/F10_089.lean}{source}\\
truncated12/catalog-chunks/F10\_090.lean & \href{https://github.com/alectionary/vibe-math/blob/master/formalization/distance-preservation/truncated12/catalog-chunks/F10_090.lean}{source} & \href{https://github.com/alectionary/vibe-math/blob/3ae21c0719cedbd928f4f7882f9b9e266567605e/formalization/distance-preservation/truncated12/catalog-chunks/F10_090.lean}{source}\\
truncated12/catalog-chunks/F10\_091.lean & \href{https://github.com/alectionary/vibe-math/blob/master/formalization/distance-preservation/truncated12/catalog-chunks/F10_091.lean}{source} & \href{https://github.com/alectionary/vibe-math/blob/3ae21c0719cedbd928f4f7882f9b9e266567605e/formalization/distance-preservation/truncated12/catalog-chunks/F10_091.lean}{source}\\
truncated12/catalog-chunks/F10\_092.lean & \href{https://github.com/alectionary/vibe-math/blob/master/formalization/distance-preservation/truncated12/catalog-chunks/F10_092.lean}{source} & \href{https://github.com/alectionary/vibe-math/blob/3ae21c0719cedbd928f4f7882f9b9e266567605e/formalization/distance-preservation/truncated12/catalog-chunks/F10_092.lean}{source}\\
truncated12/catalog-chunks/F10\_093.lean & \href{https://github.com/alectionary/vibe-math/blob/master/formalization/distance-preservation/truncated12/catalog-chunks/F10_093.lean}{source} & \href{https://github.com/alectionary/vibe-math/blob/3ae21c0719cedbd928f4f7882f9b9e266567605e/formalization/distance-preservation/truncated12/catalog-chunks/F10_093.lean}{source}\\
truncated12/catalog-chunks/F10\_094.lean & \href{https://github.com/alectionary/vibe-math/blob/master/formalization/distance-preservation/truncated12/catalog-chunks/F10_094.lean}{source} & \href{https://github.com/alectionary/vibe-math/blob/3ae21c0719cedbd928f4f7882f9b9e266567605e/formalization/distance-preservation/truncated12/catalog-chunks/F10_094.lean}{source}\\
truncated12/catalog-chunks/F10\_095.lean & \href{https://github.com/alectionary/vibe-math/blob/master/formalization/distance-preservation/truncated12/catalog-chunks/F10_095.lean}{source} & \href{https://github.com/alectionary/vibe-math/blob/3ae21c0719cedbd928f4f7882f9b9e266567605e/formalization/distance-preservation/truncated12/catalog-chunks/F10_095.lean}{source}\\
truncated12/catalog-chunks/F10\_096.lean & \href{https://github.com/alectionary/vibe-math/blob/master/formalization/distance-preservation/truncated12/catalog-chunks/F10_096.lean}{source} & \href{https://github.com/alectionary/vibe-math/blob/3ae21c0719cedbd928f4f7882f9b9e266567605e/formalization/distance-preservation/truncated12/catalog-chunks/F10_096.lean}{source}\\
truncated12/catalog-chunks/F10\_097.lean & \href{https://github.com/alectionary/vibe-math/blob/master/formalization/distance-preservation/truncated12/catalog-chunks/F10_097.lean}{source} & \href{https://github.com/alectionary/vibe-math/blob/3ae21c0719cedbd928f4f7882f9b9e266567605e/formalization/distance-preservation/truncated12/catalog-chunks/F10_097.lean}{source}\\
truncated12/catalog-chunks/F10\_098.lean & \href{https://github.com/alectionary/vibe-math/blob/master/formalization/distance-preservation/truncated12/catalog-chunks/F10_098.lean}{source} & \href{https://github.com/alectionary/vibe-math/blob/3ae21c0719cedbd928f4f7882f9b9e266567605e/formalization/distance-preservation/truncated12/catalog-chunks/F10_098.lean}{source}\\
truncated12/catalog-chunks/F10\_099.lean & \href{https://github.com/alectionary/vibe-math/blob/master/formalization/distance-preservation/truncated12/catalog-chunks/F10_099.lean}{source} & \href{https://github.com/alectionary/vibe-math/blob/3ae21c0719cedbd928f4f7882f9b9e266567605e/formalization/distance-preservation/truncated12/catalog-chunks/F10_099.lean}{source}\\
truncated12/catalog-chunks/F10\_100.lean & \href{https://github.com/alectionary/vibe-math/blob/master/formalization/distance-preservation/truncated12/catalog-chunks/F10_100.lean}{source} & \href{https://github.com/alectionary/vibe-math/blob/3ae21c0719cedbd928f4f7882f9b9e266567605e/formalization/distance-preservation/truncated12/catalog-chunks/F10_100.lean}{source}\\
truncated12/catalog-chunks/F10\_101.lean & \href{https://github.com/alectionary/vibe-math/blob/master/formalization/distance-preservation/truncated12/catalog-chunks/F10_101.lean}{source} & \href{https://github.com/alectionary/vibe-math/blob/3ae21c0719cedbd928f4f7882f9b9e266567605e/formalization/distance-preservation/truncated12/catalog-chunks/F10_101.lean}{source}\\
truncated12/catalog-chunks/F10\_102.lean & \href{https://github.com/alectionary/vibe-math/blob/master/formalization/distance-preservation/truncated12/catalog-chunks/F10_102.lean}{source} & \href{https://github.com/alectionary/vibe-math/blob/3ae21c0719cedbd928f4f7882f9b9e266567605e/formalization/distance-preservation/truncated12/catalog-chunks/F10_102.lean}{source}\\
truncated12/catalog-chunks/F10\_103.lean & \href{https://github.com/alectionary/vibe-math/blob/master/formalization/distance-preservation/truncated12/catalog-chunks/F10_103.lean}{source} & \href{https://github.com/alectionary/vibe-math/blob/3ae21c0719cedbd928f4f7882f9b9e266567605e/formalization/distance-preservation/truncated12/catalog-chunks/F10_103.lean}{source}\\
truncated12/catalog-chunks/F10\_104.lean & \href{https://github.com/alectionary/vibe-math/blob/master/formalization/distance-preservation/truncated12/catalog-chunks/F10_104.lean}{source} & \href{https://github.com/alectionary/vibe-math/blob/3ae21c0719cedbd928f4f7882f9b9e266567605e/formalization/distance-preservation/truncated12/catalog-chunks/F10_104.lean}{source}\\
truncated12/catalog-chunks/F10\_105.lean & \href{https://github.com/alectionary/vibe-math/blob/master/formalization/distance-preservation/truncated12/catalog-chunks/F10_105.lean}{source} & \href{https://github.com/alectionary/vibe-math/blob/3ae21c0719cedbd928f4f7882f9b9e266567605e/formalization/distance-preservation/truncated12/catalog-chunks/F10_105.lean}{source}\\
truncated12/catalog-chunks/F10\_106.lean & \href{https://github.com/alectionary/vibe-math/blob/master/formalization/distance-preservation/truncated12/catalog-chunks/F10_106.lean}{source} & \href{https://github.com/alectionary/vibe-math/blob/3ae21c0719cedbd928f4f7882f9b9e266567605e/formalization/distance-preservation/truncated12/catalog-chunks/F10_106.lean}{source}\\
truncated12/catalog-chunks/F10\_107.lean & \href{https://github.com/alectionary/vibe-math/blob/master/formalization/distance-preservation/truncated12/catalog-chunks/F10_107.lean}{source} & \href{https://github.com/alectionary/vibe-math/blob/3ae21c0719cedbd928f4f7882f9b9e266567605e/formalization/distance-preservation/truncated12/catalog-chunks/F10_107.lean}{source}\\
truncated12/catalog-chunks/F10\_108.lean & \href{https://github.com/alectionary/vibe-math/blob/master/formalization/distance-preservation/truncated12/catalog-chunks/F10_108.lean}{source} & \href{https://github.com/alectionary/vibe-math/blob/3ae21c0719cedbd928f4f7882f9b9e266567605e/formalization/distance-preservation/truncated12/catalog-chunks/F10_108.lean}{source}\\
truncated12/catalog-chunks/F10\_109.lean & \href{https://github.com/alectionary/vibe-math/blob/master/formalization/distance-preservation/truncated12/catalog-chunks/F10_109.lean}{source} & \href{https://github.com/alectionary/vibe-math/blob/3ae21c0719cedbd928f4f7882f9b9e266567605e/formalization/distance-preservation/truncated12/catalog-chunks/F10_109.lean}{source}\\
truncated12/catalog-chunks/F10\_110.lean & \href{https://github.com/alectionary/vibe-math/blob/master/formalization/distance-preservation/truncated12/catalog-chunks/F10_110.lean}{source} & \href{https://github.com/alectionary/vibe-math/blob/3ae21c0719cedbd928f4f7882f9b9e266567605e/formalization/distance-preservation/truncated12/catalog-chunks/F10_110.lean}{source}\\
truncated12/catalog-chunks/F10\_111.lean & \href{https://github.com/alectionary/vibe-math/blob/master/formalization/distance-preservation/truncated12/catalog-chunks/F10_111.lean}{source} & \href{https://github.com/alectionary/vibe-math/blob/3ae21c0719cedbd928f4f7882f9b9e266567605e/formalization/distance-preservation/truncated12/catalog-chunks/F10_111.lean}{source}\\
truncated12/catalog-chunks/F10\_112.lean & \href{https://github.com/alectionary/vibe-math/blob/master/formalization/distance-preservation/truncated12/catalog-chunks/F10_112.lean}{source} & \href{https://github.com/alectionary/vibe-math/blob/3ae21c0719cedbd928f4f7882f9b9e266567605e/formalization/distance-preservation/truncated12/catalog-chunks/F10_112.lean}{source}\\
truncated12/catalog-chunks/F10\_113.lean & \href{https://github.com/alectionary/vibe-math/blob/master/formalization/distance-preservation/truncated12/catalog-chunks/F10_113.lean}{source} & \href{https://github.com/alectionary/vibe-math/blob/3ae21c0719cedbd928f4f7882f9b9e266567605e/formalization/distance-preservation/truncated12/catalog-chunks/F10_113.lean}{source}\\
truncated12/catalog-chunks/F10\_114.lean & \href{https://github.com/alectionary/vibe-math/blob/master/formalization/distance-preservation/truncated12/catalog-chunks/F10_114.lean}{source} & \href{https://github.com/alectionary/vibe-math/blob/3ae21c0719cedbd928f4f7882f9b9e266567605e/formalization/distance-preservation/truncated12/catalog-chunks/F10_114.lean}{source}\\
truncated12/catalog-chunks/F10\_115.lean & \href{https://github.com/alectionary/vibe-math/blob/master/formalization/distance-preservation/truncated12/catalog-chunks/F10_115.lean}{source} & \href{https://github.com/alectionary/vibe-math/blob/3ae21c0719cedbd928f4f7882f9b9e266567605e/formalization/distance-preservation/truncated12/catalog-chunks/F10_115.lean}{source}\\
truncated12/catalog-chunks/F10\_116.lean & \href{https://github.com/alectionary/vibe-math/blob/master/formalization/distance-preservation/truncated12/catalog-chunks/F10_116.lean}{source} & \href{https://github.com/alectionary/vibe-math/blob/3ae21c0719cedbd928f4f7882f9b9e266567605e/formalization/distance-preservation/truncated12/catalog-chunks/F10_116.lean}{source}\\
truncated12/catalog-chunks/F10\_117.lean & \href{https://github.com/alectionary/vibe-math/blob/master/formalization/distance-preservation/truncated12/catalog-chunks/F10_117.lean}{source} & \href{https://github.com/alectionary/vibe-math/blob/3ae21c0719cedbd928f4f7882f9b9e266567605e/formalization/distance-preservation/truncated12/catalog-chunks/F10_117.lean}{source}\\
truncated12/catalog-chunks/F10\_118.lean & \href{https://github.com/alectionary/vibe-math/blob/master/formalization/distance-preservation/truncated12/catalog-chunks/F10_118.lean}{source} & \href{https://github.com/alectionary/vibe-math/blob/3ae21c0719cedbd928f4f7882f9b9e266567605e/formalization/distance-preservation/truncated12/catalog-chunks/F10_118.lean}{source}\\
truncated12/catalog-chunks/F10\_119.lean & \href{https://github.com/alectionary/vibe-math/blob/master/formalization/distance-preservation/truncated12/catalog-chunks/F10_119.lean}{source} & \href{https://github.com/alectionary/vibe-math/blob/3ae21c0719cedbd928f4f7882f9b9e266567605e/formalization/distance-preservation/truncated12/catalog-chunks/F10_119.lean}{source}\\
truncated12/catalog-chunks/F10\_120.lean & \href{https://github.com/alectionary/vibe-math/blob/master/formalization/distance-preservation/truncated12/catalog-chunks/F10_120.lean}{source} & \href{https://github.com/alectionary/vibe-math/blob/3ae21c0719cedbd928f4f7882f9b9e266567605e/formalization/distance-preservation/truncated12/catalog-chunks/F10_120.lean}{source}\\
truncated12/catalog-chunks/F10\_121.lean & \href{https://github.com/alectionary/vibe-math/blob/master/formalization/distance-preservation/truncated12/catalog-chunks/F10_121.lean}{source} & \href{https://github.com/alectionary/vibe-math/blob/3ae21c0719cedbd928f4f7882f9b9e266567605e/formalization/distance-preservation/truncated12/catalog-chunks/F10_121.lean}{source}\\
truncated12/catalog-chunks/F10\_122.lean & \href{https://github.com/alectionary/vibe-math/blob/master/formalization/distance-preservation/truncated12/catalog-chunks/F10_122.lean}{source} & \href{https://github.com/alectionary/vibe-math/blob/3ae21c0719cedbd928f4f7882f9b9e266567605e/formalization/distance-preservation/truncated12/catalog-chunks/F10_122.lean}{source}\\
truncated12/catalog-chunks/F10\_123.lean & \href{https://github.com/alectionary/vibe-math/blob/master/formalization/distance-preservation/truncated12/catalog-chunks/F10_123.lean}{source} & \href{https://github.com/alectionary/vibe-math/blob/3ae21c0719cedbd928f4f7882f9b9e266567605e/formalization/distance-preservation/truncated12/catalog-chunks/F10_123.lean}{source}\\
truncated12/catalog-chunks/F10\_124.lean & \href{https://github.com/alectionary/vibe-math/blob/master/formalization/distance-preservation/truncated12/catalog-chunks/F10_124.lean}{source} & \href{https://github.com/alectionary/vibe-math/blob/3ae21c0719cedbd928f4f7882f9b9e266567605e/formalization/distance-preservation/truncated12/catalog-chunks/F10_124.lean}{source}\\
truncated12/Catalog10.lean & \href{https://github.com/alectionary/vibe-math/blob/master/formalization/distance-preservation/truncated12/Catalog10.lean}{source} & \href{https://github.com/alectionary/vibe-math/blob/3ae21c0719cedbd928f4f7882f9b9e266567605e/formalization/distance-preservation/truncated12/Catalog10.lean}{source}\\
truncated12/CoverageCore.lean & \href{https://github.com/alectionary/vibe-math/blob/master/formalization/distance-preservation/truncated12/CoverageCore.lean}{source} & \href{https://github.com/alectionary/vibe-math/blob/3ae21c0719cedbd928f4f7882f9b9e266567605e/formalization/distance-preservation/truncated12/CoverageCore.lean}{source}\\
truncated12/Coverage10.lean & \href{https://github.com/alectionary/vibe-math/blob/master/formalization/distance-preservation/truncated12/Coverage10.lean}{source} & \href{https://github.com/alectionary/vibe-math/blob/3ae21c0719cedbd928f4f7882f9b9e266567605e/formalization/distance-preservation/truncated12/Coverage10.lean}{source}\\
FamilyWords.lean & \href{https://github.com/alectionary/vibe-math/blob/master/formalization/distance-preservation/FamilyWords.lean}{source} & \href{https://github.com/alectionary/vibe-math/blob/3ae21c0719cedbd928f4f7882f9b9e266567605e/formalization/distance-preservation/FamilyWords.lean}{source}\\
CoverWeight.lean & \href{https://github.com/alectionary/vibe-math/blob/master/formalization/distance-preservation/CoverWeight.lean}{source} & \href{https://github.com/alectionary/vibe-math/blob/3ae21c0719cedbd928f4f7882f9b9e266567605e/formalization/distance-preservation/CoverWeight.lean}{source}\\
SourceModel.lean & \href{https://github.com/alectionary/vibe-math/blob/master/formalization/distance-preservation/SourceModel.lean}{source} & \href{https://github.com/alectionary/vibe-math/blob/3ae21c0719cedbd928f4f7882f9b9e266567605e/formalization/distance-preservation/SourceModel.lean}{source}\\
GeneralLP.lean & \href{https://github.com/alectionary/vibe-math/blob/master/formalization/distance-preservation/GeneralLP.lean}{source} & \href{https://github.com/alectionary/vibe-math/blob/3ae21c0719cedbd928f4f7882f9b9e266567605e/formalization/distance-preservation/GeneralLP.lean}{source}\\
SmallStopping.lean & \href{https://github.com/alectionary/vibe-math/blob/master/formalization/distance-preservation/SmallStopping.lean}{source} & \href{https://github.com/alectionary/vibe-math/blob/3ae21c0719cedbd928f4f7882f9b9e266567605e/formalization/distance-preservation/SmallStopping.lean}{source}\\
CyclicCover.lean & \href{https://github.com/alectionary/vibe-math/blob/master/formalization/distance-preservation/CyclicCover.lean}{source} & \href{https://github.com/alectionary/vibe-math/blob/3ae21c0719cedbd928f4f7882f9b9e266567605e/formalization/distance-preservation/CyclicCover.lean}{source}\\
GeneralLPPhysical.lean & \href{https://github.com/alectionary/vibe-math/blob/master/formalization/distance-preservation/GeneralLPPhysical.lean}{source} & \href{https://github.com/alectionary/vibe-math/blob/3ae21c0719cedbd928f4f7882f9b9e266567605e/formalization/distance-preservation/GeneralLPPhysical.lean}{source}\\
TannerSix.lean & \href{https://github.com/alectionary/vibe-math/blob/master/formalization/distance-preservation/TannerSix.lean}{source} & \href{https://github.com/alectionary/vibe-math/blob/3ae21c0719cedbd928f4f7882f9b9e266567605e/formalization/distance-preservation/TannerSix.lean}{source}\\
LowWeightChecks.lean & \href{https://github.com/alectionary/vibe-math/blob/master/formalization/distance-preservation/LowWeightChecks.lean}{source} & \href{https://github.com/alectionary/vibe-math/blob/3ae21c0719cedbd928f4f7882f9b9e266567605e/formalization/distance-preservation/LowWeightChecks.lean}{source}\\
truncated12/ActualIncidence.lean & \href{https://github.com/alectionary/vibe-math/blob/master/formalization/distance-preservation/truncated12/ActualIncidence.lean}{source} & \href{https://github.com/alectionary/vibe-math/blob/3ae21c0719cedbd928f4f7882f9b9e266567605e/formalization/distance-preservation/truncated12/ActualIncidence.lean}{source}\\
truncated12/NormalGraph.lean & \href{https://github.com/alectionary/vibe-math/blob/master/formalization/distance-preservation/truncated12/NormalGraph.lean}{source} & \href{https://github.com/alectionary/vibe-math/blob/3ae21c0719cedbd928f4f7882f9b9e266567605e/formalization/distance-preservation/truncated12/NormalGraph.lean}{source}\\
truncated12/BlockParity.lean & \href{https://github.com/alectionary/vibe-math/blob/master/formalization/distance-preservation/truncated12/BlockParity.lean}{source} & \href{https://github.com/alectionary/vibe-math/blob/3ae21c0719cedbd928f4f7882f9b9e266567605e/formalization/distance-preservation/truncated12/BlockParity.lean}{source}\\
truncated12/BlockCounts.lean & \href{https://github.com/alectionary/vibe-math/blob/master/formalization/distance-preservation/truncated12/BlockCounts.lean}{source} & \href{https://github.com/alectionary/vibe-math/blob/3ae21c0719cedbd928f4f7882f9b9e266567605e/formalization/distance-preservation/truncated12/BlockCounts.lean}{source}\\
truncated12/SourceNormal.lean & \href{https://github.com/alectionary/vibe-math/blob/master/formalization/distance-preservation/truncated12/SourceNormal.lean}{source} & \href{https://github.com/alectionary/vibe-math/blob/3ae21c0719cedbd928f4f7882f9b9e266567605e/formalization/distance-preservation/truncated12/SourceNormal.lean}{source}\\
truncated12/BucketCounts.lean & \href{https://github.com/alectionary/vibe-math/blob/master/formalization/distance-preservation/truncated12/BucketCounts.lean}{source} & \href{https://github.com/alectionary/vibe-math/blob/3ae21c0719cedbd928f4f7882f9b9e266567605e/formalization/distance-preservation/truncated12/BucketCounts.lean}{source}\\
truncated12/TorsionClassical.lean & \href{https://github.com/alectionary/vibe-math/blob/master/formalization/distance-preservation/truncated12/TorsionClassical.lean}{source} & \href{https://github.com/alectionary/vibe-math/blob/3ae21c0719cedbd928f4f7882f9b9e266567605e/formalization/distance-preservation/truncated12/TorsionClassical.lean}{source}\\
truncated12/IntegerSoundness.lean & \href{https://github.com/alectionary/vibe-math/blob/master/formalization/distance-preservation/truncated12/IntegerSoundness.lean}{source} & \href{https://github.com/alectionary/vibe-math/blob/3ae21c0719cedbd928f4f7882f9b9e266567605e/formalization/distance-preservation/truncated12/IntegerSoundness.lean}{source}\\
truncated12/SingletonLabels.lean & \href{https://github.com/alectionary/vibe-math/blob/master/formalization/distance-preservation/truncated12/SingletonLabels.lean}{source} & \href{https://github.com/alectionary/vibe-math/blob/3ae21c0719cedbd928f4f7882f9b9e266567605e/formalization/distance-preservation/truncated12/SingletonLabels.lean}{source}\\
truncated12/SupportSingletons.lean & \href{https://github.com/alectionary/vibe-math/blob/master/formalization/distance-preservation/truncated12/SupportSingletons.lean}{source} & \href{https://github.com/alectionary/vibe-math/blob/3ae21c0719cedbd928f4f7882f9b9e266567605e/formalization/distance-preservation/truncated12/SupportSingletons.lean}{source}\\
truncated12/SupportPairs.lean & \href{https://github.com/alectionary/vibe-math/blob/master/formalization/distance-preservation/truncated12/SupportPairs.lean}{source} & \href{https://github.com/alectionary/vibe-math/blob/3ae21c0719cedbd928f4f7882f9b9e266567605e/formalization/distance-preservation/truncated12/SupportPairs.lean}{source}\\
truncated12/PairMatching.lean & \href{https://github.com/alectionary/vibe-math/blob/master/formalization/distance-preservation/truncated12/PairMatching.lean}{source} & \href{https://github.com/alectionary/vibe-math/blob/3ae21c0719cedbd928f4f7882f9b9e266567605e/formalization/distance-preservation/truncated12/PairMatching.lean}{source}\\
truncated12/KernelCatalog.lean & \href{https://github.com/alectionary/vibe-math/blob/master/formalization/distance-preservation/truncated12/KernelCatalog.lean}{source} & \href{https://github.com/alectionary/vibe-math/blob/3ae21c0719cedbd928f4f7882f9b9e266567605e/formalization/distance-preservation/truncated12/KernelCatalog.lean}{source}\\
truncated12/MatchingRows.lean & \href{https://github.com/alectionary/vibe-math/blob/master/formalization/distance-preservation/truncated12/MatchingRows.lean}{source} & \href{https://github.com/alectionary/vibe-math/blob/3ae21c0719cedbd928f4f7882f9b9e266567605e/formalization/distance-preservation/truncated12/MatchingRows.lean}{source}\\
truncated12/GaugeTransport.lean & \href{https://github.com/alectionary/vibe-math/blob/master/formalization/distance-preservation/truncated12/GaugeTransport.lean}{source} & \href{https://github.com/alectionary/vibe-math/blob/3ae21c0719cedbd928f4f7882f9b9e266567605e/formalization/distance-preservation/truncated12/GaugeTransport.lean}{source}\\
truncated12/Normalization.lean & \href{https://github.com/alectionary/vibe-math/blob/master/formalization/distance-preservation/truncated12/Normalization.lean}{source} & \href{https://github.com/alectionary/vibe-math/blob/3ae21c0719cedbd928f4f7882f9b9e266567605e/formalization/distance-preservation/truncated12/Normalization.lean}{source}\\
truncated12/SourceCycleFacts.lean & \href{https://github.com/alectionary/vibe-math/blob/master/formalization/distance-preservation/truncated12/SourceCycleFacts.lean}{source} & \href{https://github.com/alectionary/vibe-math/blob/3ae21c0719cedbd928f4f7882f9b9e266567605e/formalization/distance-preservation/truncated12/SourceCycleFacts.lean}{source}\\
truncated12/PureDistance.lean & \href{https://github.com/alectionary/vibe-math/blob/master/formalization/distance-preservation/truncated12/PureDistance.lean}{source} & \href{https://github.com/alectionary/vibe-math/blob/3ae21c0719cedbd928f4f7882f9b9e266567605e/formalization/distance-preservation/truncated12/PureDistance.lean}{source}\\
truncated12/CatalogTables.lean & \href{https://github.com/alectionary/vibe-math/blob/master/formalization/distance-preservation/truncated12/CatalogTables.lean}{source} & \href{https://github.com/alectionary/vibe-math/blob/3ae21c0719cedbd928f4f7882f9b9e266567605e/formalization/distance-preservation/truncated12/CatalogTables.lean}{source}\\
truncated12/PackedTablesAudit.lean & \href{https://github.com/alectionary/vibe-math/blob/master/formalization/distance-preservation/truncated12/PackedTablesAudit.lean}{source} & \href{https://github.com/alectionary/vibe-math/blob/3ae21c0719cedbd928f4f7882f9b9e266567605e/formalization/distance-preservation/truncated12/PackedTablesAudit.lean}{source}\\
truncated12/TargetObstructions.lean & \href{https://github.com/alectionary/vibe-math/blob/master/formalization/distance-preservation/truncated12/TargetObstructions.lean}{source} & \href{https://github.com/alectionary/vibe-math/blob/3ae21c0719cedbd928f4f7882f9b9e266567605e/formalization/distance-preservation/truncated12/TargetObstructions.lean}{source}\\
truncated12/PatternRealizations.lean & \href{https://github.com/alectionary/vibe-math/blob/master/formalization/distance-preservation/truncated12/PatternRealizations.lean}{source} & \href{https://github.com/alectionary/vibe-math/blob/3ae21c0719cedbd928f4f7882f9b9e266567605e/formalization/distance-preservation/truncated12/PatternRealizations.lean}{source}\\
truncated12/GeneralSymmetry.lean & \href{https://github.com/alectionary/vibe-math/blob/master/formalization/distance-preservation/truncated12/GeneralSymmetry.lean}{source} & \href{https://github.com/alectionary/vibe-math/blob/3ae21c0719cedbd928f4f7882f9b9e266567605e/formalization/distance-preservation/truncated12/GeneralSymmetry.lean}{source}\\
truncated12/SpanCoefficients.lean & \href{https://github.com/alectionary/vibe-math/blob/master/formalization/distance-preservation/truncated12/SpanCoefficients.lean}{source} & \href{https://github.com/alectionary/vibe-math/blob/3ae21c0719cedbd928f4f7882f9b9e266567605e/formalization/distance-preservation/truncated12/SpanCoefficients.lean}{source}\\
truncated12/RowCancellation.lean & \href{https://github.com/alectionary/vibe-math/blob/master/formalization/distance-preservation/truncated12/RowCancellation.lean}{source} & \href{https://github.com/alectionary/vibe-math/blob/3ae21c0719cedbd928f4f7882f9b9e266567605e/formalization/distance-preservation/truncated12/RowCancellation.lean}{source}\\
truncated12/ConstantRows.lean & \href{https://github.com/alectionary/vibe-math/blob/master/formalization/distance-preservation/truncated12/ConstantRows.lean}{source} & \href{https://github.com/alectionary/vibe-math/blob/3ae21c0719cedbd928f4f7882f9b9e266567605e/formalization/distance-preservation/truncated12/ConstantRows.lean}{source}\\
truncated12/FiniteLabels.lean & \href{https://github.com/alectionary/vibe-math/blob/master/formalization/distance-preservation/truncated12/FiniteLabels.lean}{source} & \href{https://github.com/alectionary/vibe-math/blob/3ae21c0719cedbd928f4f7882f9b9e266567605e/formalization/distance-preservation/truncated12/FiniteLabels.lean}{source}\\
truncated12/catalog-chunks/F7\_000.lean & \href{https://github.com/alectionary/vibe-math/blob/master/formalization/distance-preservation/truncated12/catalog-chunks/F7_000.lean}{source} & \href{https://github.com/alectionary/vibe-math/blob/3ae21c0719cedbd928f4f7882f9b9e266567605e/formalization/distance-preservation/truncated12/catalog-chunks/F7_000.lean}{source}\\
truncated12/Catalog7.lean & \href{https://github.com/alectionary/vibe-math/blob/master/formalization/distance-preservation/truncated12/Catalog7.lean}{source} & \href{https://github.com/alectionary/vibe-math/blob/3ae21c0719cedbd928f4f7882f9b9e266567605e/formalization/distance-preservation/truncated12/Catalog7.lean}{source}\\
truncated12/SourceSeven.lean & \href{https://github.com/alectionary/vibe-math/blob/master/formalization/distance-preservation/truncated12/SourceSeven.lean}{source} & \href{https://github.com/alectionary/vibe-math/blob/3ae21c0719cedbd928f4f7882f9b9e266567605e/formalization/distance-preservation/truncated12/SourceSeven.lean}{source}\\
truncated12/SocketEdges.lean & \href{https://github.com/alectionary/vibe-math/blob/master/formalization/distance-preservation/truncated12/SocketEdges.lean}{source} & \href{https://github.com/alectionary/vibe-math/blob/3ae21c0719cedbd928f4f7882f9b9e266567605e/formalization/distance-preservation/truncated12/SocketEdges.lean}{source}\\
truncated12/SocketArithmetic.lean & \href{https://github.com/alectionary/vibe-math/blob/master/formalization/distance-preservation/truncated12/SocketArithmetic.lean}{source} & \href{https://github.com/alectionary/vibe-math/blob/3ae21c0719cedbd928f4f7882f9b9e266567605e/formalization/distance-preservation/truncated12/SocketArithmetic.lean}{source}\\
truncated12/BucketSums.lean & \href{https://github.com/alectionary/vibe-math/blob/master/formalization/distance-preservation/truncated12/BucketSums.lean}{source} & \href{https://github.com/alectionary/vibe-math/blob/3ae21c0719cedbd928f4f7882f9b9e266567605e/formalization/distance-preservation/truncated12/BucketSums.lean}{source}\\
truncated12/LocalMantel.lean & \href{https://github.com/alectionary/vibe-math/blob/master/formalization/distance-preservation/truncated12/LocalMantel.lean}{source} & \href{https://github.com/alectionary/vibe-math/blob/3ae21c0719cedbd928f4f7882f9b9e266567605e/formalization/distance-preservation/truncated12/LocalMantel.lean}{source}\\
truncated12/CountPatterns.lean & \href{https://github.com/alectionary/vibe-math/blob/master/formalization/distance-preservation/truncated12/CountPatterns.lean}{source} & \href{https://github.com/alectionary/vibe-math/blob/3ae21c0719cedbd928f4f7882f9b9e266567605e/formalization/distance-preservation/truncated12/CountPatterns.lean}{source}\\
truncated12/SourceProfiles.lean & \href{https://github.com/alectionary/vibe-math/blob/master/formalization/distance-preservation/truncated12/SourceProfiles.lean}{source} & \href{https://github.com/alectionary/vibe-math/blob/3ae21c0719cedbd928f4f7882f9b9e266567605e/formalization/distance-preservation/truncated12/SourceProfiles.lean}{source}\\
truncated12/PauliBridge.lean & \href{https://github.com/alectionary/vibe-math/blob/master/formalization/distance-preservation/truncated12/PauliBridge.lean}{source} & \href{https://github.com/alectionary/vibe-math/blob/3ae21c0719cedbd928f4f7882f9b9e266567605e/formalization/distance-preservation/truncated12/PauliBridge.lean}{source}\\
truncated12/ThroughNine.lean & \href{https://github.com/alectionary/vibe-math/blob/master/formalization/distance-preservation/truncated12/ThroughNine.lean}{source} & \href{https://github.com/alectionary/vibe-math/blob/3ae21c0719cedbd928f4f7882f9b9e266567605e/formalization/distance-preservation/truncated12/ThroughNine.lean}{source}\\
truncated12/TenGeometry.lean & \href{https://github.com/alectionary/vibe-math/blob/master/formalization/distance-preservation/truncated12/TenGeometry.lean}{source} & \href{https://github.com/alectionary/vibe-math/blob/3ae21c0719cedbd928f4f7882f9b9e266567605e/formalization/distance-preservation/truncated12/TenGeometry.lean}{source}\\
truncated12/TenIndex.lean & \href{https://github.com/alectionary/vibe-math/blob/master/formalization/distance-preservation/truncated12/TenIndex.lean}{source} & \href{https://github.com/alectionary/vibe-math/blob/3ae21c0719cedbd928f4f7882f9b9e266567605e/formalization/distance-preservation/truncated12/TenIndex.lean}{source}\\
truncated12/TenRealization.lean & \href{https://github.com/alectionary/vibe-math/blob/master/formalization/distance-preservation/truncated12/TenRealization.lean}{source} & \href{https://github.com/alectionary/vibe-math/blob/3ae21c0719cedbd928f4f7882f9b9e266567605e/formalization/distance-preservation/truncated12/TenRealization.lean}{source}\\
truncated12/TenEncoding.lean & \href{https://github.com/alectionary/vibe-math/blob/master/formalization/distance-preservation/truncated12/TenEncoding.lean}{source} & \href{https://github.com/alectionary/vibe-math/blob/3ae21c0719cedbd928f4f7882f9b9e266567605e/formalization/distance-preservation/truncated12/TenEncoding.lean}{source}\\
truncated12/SourceTen.lean & \href{https://github.com/alectionary/vibe-math/blob/master/formalization/distance-preservation/truncated12/SourceTen.lean}{source} & \href{https://github.com/alectionary/vibe-math/blob/3ae21c0719cedbd928f4f7882f9b9e266567605e/formalization/distance-preservation/truncated12/SourceTen.lean}{source}\\
truncated12/ExceptionalGraph.lean & \href{https://github.com/alectionary/vibe-math/blob/master/formalization/distance-preservation/truncated12/ExceptionalGraph.lean}{source} & \href{https://github.com/alectionary/vibe-math/blob/3ae21c0719cedbd928f4f7882f9b9e266567605e/formalization/distance-preservation/truncated12/ExceptionalGraph.lean}{source}\\
truncated12/SourcePartners.lean & \href{https://github.com/alectionary/vibe-math/blob/master/formalization/distance-preservation/truncated12/SourcePartners.lean}{source} & \href{https://github.com/alectionary/vibe-math/blob/3ae21c0719cedbd928f4f7882f9b9e266567605e/formalization/distance-preservation/truncated12/SourcePartners.lean}{source}\\
truncated12/K23Numbering.lean & \href{https://github.com/alectionary/vibe-math/blob/master/formalization/distance-preservation/truncated12/K23Numbering.lean}{source} & \href{https://github.com/alectionary/vibe-math/blob/3ae21c0719cedbd928f4f7882f9b9e266567605e/formalization/distance-preservation/truncated12/K23Numbering.lean}{source}\\
truncated12/K23Colors.lean & \href{https://github.com/alectionary/vibe-math/blob/master/formalization/distance-preservation/truncated12/K23Colors.lean}{source} & \href{https://github.com/alectionary/vibe-math/blob/3ae21c0719cedbd928f4f7882f9b9e266567605e/formalization/distance-preservation/truncated12/K23Colors.lean}{source}\\
truncated12/K23MissingRows.lean & \href{https://github.com/alectionary/vibe-math/blob/master/formalization/distance-preservation/truncated12/K23MissingRows.lean}{source} & \href{https://github.com/alectionary/vibe-math/blob/3ae21c0719cedbd928f4f7882f9b9e266567605e/formalization/distance-preservation/truncated12/K23MissingRows.lean}{source}\\
truncated12/K23LeafMatches.lean & \href{https://github.com/alectionary/vibe-math/blob/master/formalization/distance-preservation/truncated12/K23LeafMatches.lean}{source} & \href{https://github.com/alectionary/vibe-math/blob/3ae21c0719cedbd928f4f7882f9b9e266567605e/formalization/distance-preservation/truncated12/K23LeafMatches.lean}{source}\\
truncated12/ElevenGeometry.lean & \href{https://github.com/alectionary/vibe-math/blob/master/formalization/distance-preservation/truncated12/ElevenGeometry.lean}{source} & \href{https://github.com/alectionary/vibe-math/blob/3ae21c0719cedbd928f4f7882f9b9e266567605e/formalization/distance-preservation/truncated12/ElevenGeometry.lean}{source}\\
truncated12/ElevenIndex.lean & \href{https://github.com/alectionary/vibe-math/blob/master/formalization/distance-preservation/truncated12/ElevenIndex.lean}{source} & \href{https://github.com/alectionary/vibe-math/blob/3ae21c0719cedbd928f4f7882f9b9e266567605e/formalization/distance-preservation/truncated12/ElevenIndex.lean}{source}\\
truncated12/ElevenRealization.lean & \href{https://github.com/alectionary/vibe-math/blob/master/formalization/distance-preservation/truncated12/ElevenRealization.lean}{source} & \href{https://github.com/alectionary/vibe-math/blob/3ae21c0719cedbd928f4f7882f9b9e266567605e/formalization/distance-preservation/truncated12/ElevenRealization.lean}{source}\\
truncated12/ElevenEncoding.lean & \href{https://github.com/alectionary/vibe-math/blob/master/formalization/distance-preservation/truncated12/ElevenEncoding.lean}{source} & \href{https://github.com/alectionary/vibe-math/blob/3ae21c0719cedbd928f4f7882f9b9e266567605e/formalization/distance-preservation/truncated12/ElevenEncoding.lean}{source}\\
truncated12/catalog-chunks/F11\_000.lean & \href{https://github.com/alectionary/vibe-math/blob/master/formalization/distance-preservation/truncated12/catalog-chunks/F11_000.lean}{source} & \href{https://github.com/alectionary/vibe-math/blob/3ae21c0719cedbd928f4f7882f9b9e266567605e/formalization/distance-preservation/truncated12/catalog-chunks/F11_000.lean}{source}\\
truncated12/catalog-chunks/F11\_001.lean & \href{https://github.com/alectionary/vibe-math/blob/master/formalization/distance-preservation/truncated12/catalog-chunks/F11_001.lean}{source} & \href{https://github.com/alectionary/vibe-math/blob/3ae21c0719cedbd928f4f7882f9b9e266567605e/formalization/distance-preservation/truncated12/catalog-chunks/F11_001.lean}{source}\\
truncated12/catalog-chunks/F11\_002.lean & \href{https://github.com/alectionary/vibe-math/blob/master/formalization/distance-preservation/truncated12/catalog-chunks/F11_002.lean}{source} & \href{https://github.com/alectionary/vibe-math/blob/3ae21c0719cedbd928f4f7882f9b9e266567605e/formalization/distance-preservation/truncated12/catalog-chunks/F11_002.lean}{source}\\
truncated12/catalog-chunks/F11\_003.lean & \href{https://github.com/alectionary/vibe-math/blob/master/formalization/distance-preservation/truncated12/catalog-chunks/F11_003.lean}{source} & \href{https://github.com/alectionary/vibe-math/blob/3ae21c0719cedbd928f4f7882f9b9e266567605e/formalization/distance-preservation/truncated12/catalog-chunks/F11_003.lean}{source}\\
truncated12/catalog-chunks/F11\_004.lean & \href{https://github.com/alectionary/vibe-math/blob/master/formalization/distance-preservation/truncated12/catalog-chunks/F11_004.lean}{source} & \href{https://github.com/alectionary/vibe-math/blob/3ae21c0719cedbd928f4f7882f9b9e266567605e/formalization/distance-preservation/truncated12/catalog-chunks/F11_004.lean}{source}\\
truncated12/catalog-chunks/F11\_005.lean & \href{https://github.com/alectionary/vibe-math/blob/master/formalization/distance-preservation/truncated12/catalog-chunks/F11_005.lean}{source} & \href{https://github.com/alectionary/vibe-math/blob/3ae21c0719cedbd928f4f7882f9b9e266567605e/formalization/distance-preservation/truncated12/catalog-chunks/F11_005.lean}{source}\\
truncated12/catalog-chunks/F11\_006.lean & \href{https://github.com/alectionary/vibe-math/blob/master/formalization/distance-preservation/truncated12/catalog-chunks/F11_006.lean}{source} & \href{https://github.com/alectionary/vibe-math/blob/3ae21c0719cedbd928f4f7882f9b9e266567605e/formalization/distance-preservation/truncated12/catalog-chunks/F11_006.lean}{source}\\
truncated12/catalog-chunks/F11\_007.lean & \href{https://github.com/alectionary/vibe-math/blob/master/formalization/distance-preservation/truncated12/catalog-chunks/F11_007.lean}{source} & \href{https://github.com/alectionary/vibe-math/blob/3ae21c0719cedbd928f4f7882f9b9e266567605e/formalization/distance-preservation/truncated12/catalog-chunks/F11_007.lean}{source}\\
truncated12/catalog-chunks/F11\_008.lean & \href{https://github.com/alectionary/vibe-math/blob/master/formalization/distance-preservation/truncated12/catalog-chunks/F11_008.lean}{source} & \href{https://github.com/alectionary/vibe-math/blob/3ae21c0719cedbd928f4f7882f9b9e266567605e/formalization/distance-preservation/truncated12/catalog-chunks/F11_008.lean}{source}\\
truncated12/catalog-chunks/F11\_009.lean & \href{https://github.com/alectionary/vibe-math/blob/master/formalization/distance-preservation/truncated12/catalog-chunks/F11_009.lean}{source} & \href{https://github.com/alectionary/vibe-math/blob/3ae21c0719cedbd928f4f7882f9b9e266567605e/formalization/distance-preservation/truncated12/catalog-chunks/F11_009.lean}{source}\\
truncated12/catalog-chunks/F11\_010.lean & \href{https://github.com/alectionary/vibe-math/blob/master/formalization/distance-preservation/truncated12/catalog-chunks/F11_010.lean}{source} & \href{https://github.com/alectionary/vibe-math/blob/3ae21c0719cedbd928f4f7882f9b9e266567605e/formalization/distance-preservation/truncated12/catalog-chunks/F11_010.lean}{source}\\
truncated12/catalog-chunks/F11\_011.lean & \href{https://github.com/alectionary/vibe-math/blob/master/formalization/distance-preservation/truncated12/catalog-chunks/F11_011.lean}{source} & \href{https://github.com/alectionary/vibe-math/blob/3ae21c0719cedbd928f4f7882f9b9e266567605e/formalization/distance-preservation/truncated12/catalog-chunks/F11_011.lean}{source}\\
truncated12/catalog-chunks/F11\_012.lean & \href{https://github.com/alectionary/vibe-math/blob/master/formalization/distance-preservation/truncated12/catalog-chunks/F11_012.lean}{source} & \href{https://github.com/alectionary/vibe-math/blob/3ae21c0719cedbd928f4f7882f9b9e266567605e/formalization/distance-preservation/truncated12/catalog-chunks/F11_012.lean}{source}\\
truncated12/catalog-chunks/F11\_013.lean & \href{https://github.com/alectionary/vibe-math/blob/master/formalization/distance-preservation/truncated12/catalog-chunks/F11_013.lean}{source} & \href{https://github.com/alectionary/vibe-math/blob/3ae21c0719cedbd928f4f7882f9b9e266567605e/formalization/distance-preservation/truncated12/catalog-chunks/F11_013.lean}{source}\\
truncated12/catalog-chunks/F11\_014.lean & \href{https://github.com/alectionary/vibe-math/blob/master/formalization/distance-preservation/truncated12/catalog-chunks/F11_014.lean}{source} & \href{https://github.com/alectionary/vibe-math/blob/3ae21c0719cedbd928f4f7882f9b9e266567605e/formalization/distance-preservation/truncated12/catalog-chunks/F11_014.lean}{source}\\
truncated12/catalog-chunks/F11\_015.lean & \href{https://github.com/alectionary/vibe-math/blob/master/formalization/distance-preservation/truncated12/catalog-chunks/F11_015.lean}{source} & \href{https://github.com/alectionary/vibe-math/blob/3ae21c0719cedbd928f4f7882f9b9e266567605e/formalization/distance-preservation/truncated12/catalog-chunks/F11_015.lean}{source}\\
truncated12/catalog-chunks/F11\_016.lean & \href{https://github.com/alectionary/vibe-math/blob/master/formalization/distance-preservation/truncated12/catalog-chunks/F11_016.lean}{source} & \href{https://github.com/alectionary/vibe-math/blob/3ae21c0719cedbd928f4f7882f9b9e266567605e/formalization/distance-preservation/truncated12/catalog-chunks/F11_016.lean}{source}\\
truncated12/catalog-chunks/F11\_017.lean & \href{https://github.com/alectionary/vibe-math/blob/master/formalization/distance-preservation/truncated12/catalog-chunks/F11_017.lean}{source} & \href{https://github.com/alectionary/vibe-math/blob/3ae21c0719cedbd928f4f7882f9b9e266567605e/formalization/distance-preservation/truncated12/catalog-chunks/F11_017.lean}{source}\\
truncated12/catalog-chunks/F11\_018.lean & \href{https://github.com/alectionary/vibe-math/blob/master/formalization/distance-preservation/truncated12/catalog-chunks/F11_018.lean}{source} & \href{https://github.com/alectionary/vibe-math/blob/3ae21c0719cedbd928f4f7882f9b9e266567605e/formalization/distance-preservation/truncated12/catalog-chunks/F11_018.lean}{source}\\
truncated12/catalog-chunks/F11\_019.lean & \href{https://github.com/alectionary/vibe-math/blob/master/formalization/distance-preservation/truncated12/catalog-chunks/F11_019.lean}{source} & \href{https://github.com/alectionary/vibe-math/blob/3ae21c0719cedbd928f4f7882f9b9e266567605e/formalization/distance-preservation/truncated12/catalog-chunks/F11_019.lean}{source}\\
truncated12/catalog-chunks/F11\_020.lean & \href{https://github.com/alectionary/vibe-math/blob/master/formalization/distance-preservation/truncated12/catalog-chunks/F11_020.lean}{source} & \href{https://github.com/alectionary/vibe-math/blob/3ae21c0719cedbd928f4f7882f9b9e266567605e/formalization/distance-preservation/truncated12/catalog-chunks/F11_020.lean}{source}\\
truncated12/catalog-chunks/F11\_021.lean & \href{https://github.com/alectionary/vibe-math/blob/master/formalization/distance-preservation/truncated12/catalog-chunks/F11_021.lean}{source} & \href{https://github.com/alectionary/vibe-math/blob/3ae21c0719cedbd928f4f7882f9b9e266567605e/formalization/distance-preservation/truncated12/catalog-chunks/F11_021.lean}{source}\\
truncated12/catalog-chunks/F11\_022.lean & \href{https://github.com/alectionary/vibe-math/blob/master/formalization/distance-preservation/truncated12/catalog-chunks/F11_022.lean}{source} & \href{https://github.com/alectionary/vibe-math/blob/3ae21c0719cedbd928f4f7882f9b9e266567605e/formalization/distance-preservation/truncated12/catalog-chunks/F11_022.lean}{source}\\
truncated12/catalog-chunks/F11\_023.lean & \href{https://github.com/alectionary/vibe-math/blob/master/formalization/distance-preservation/truncated12/catalog-chunks/F11_023.lean}{source} & \href{https://github.com/alectionary/vibe-math/blob/3ae21c0719cedbd928f4f7882f9b9e266567605e/formalization/distance-preservation/truncated12/catalog-chunks/F11_023.lean}{source}\\
truncated12/catalog-chunks/F11\_024.lean & \href{https://github.com/alectionary/vibe-math/blob/master/formalization/distance-preservation/truncated12/catalog-chunks/F11_024.lean}{source} & \href{https://github.com/alectionary/vibe-math/blob/3ae21c0719cedbd928f4f7882f9b9e266567605e/formalization/distance-preservation/truncated12/catalog-chunks/F11_024.lean}{source}\\
truncated12/catalog-chunks/F11\_025.lean & \href{https://github.com/alectionary/vibe-math/blob/master/formalization/distance-preservation/truncated12/catalog-chunks/F11_025.lean}{source} & \href{https://github.com/alectionary/vibe-math/blob/3ae21c0719cedbd928f4f7882f9b9e266567605e/formalization/distance-preservation/truncated12/catalog-chunks/F11_025.lean}{source}\\
truncated12/catalog-chunks/F11\_026.lean & \href{https://github.com/alectionary/vibe-math/blob/master/formalization/distance-preservation/truncated12/catalog-chunks/F11_026.lean}{source} & \href{https://github.com/alectionary/vibe-math/blob/3ae21c0719cedbd928f4f7882f9b9e266567605e/formalization/distance-preservation/truncated12/catalog-chunks/F11_026.lean}{source}\\
truncated12/catalog-chunks/F11\_027.lean & \href{https://github.com/alectionary/vibe-math/blob/master/formalization/distance-preservation/truncated12/catalog-chunks/F11_027.lean}{source} & \href{https://github.com/alectionary/vibe-math/blob/3ae21c0719cedbd928f4f7882f9b9e266567605e/formalization/distance-preservation/truncated12/catalog-chunks/F11_027.lean}{source}\\
truncated12/catalog-chunks/F11\_028.lean & \href{https://github.com/alectionary/vibe-math/blob/master/formalization/distance-preservation/truncated12/catalog-chunks/F11_028.lean}{source} & \href{https://github.com/alectionary/vibe-math/blob/3ae21c0719cedbd928f4f7882f9b9e266567605e/formalization/distance-preservation/truncated12/catalog-chunks/F11_028.lean}{source}\\
truncated12/catalog-chunks/F11\_029.lean & \href{https://github.com/alectionary/vibe-math/blob/master/formalization/distance-preservation/truncated12/catalog-chunks/F11_029.lean}{source} & \href{https://github.com/alectionary/vibe-math/blob/3ae21c0719cedbd928f4f7882f9b9e266567605e/formalization/distance-preservation/truncated12/catalog-chunks/F11_029.lean}{source}\\
truncated12/catalog-chunks/F11\_030.lean & \href{https://github.com/alectionary/vibe-math/blob/master/formalization/distance-preservation/truncated12/catalog-chunks/F11_030.lean}{source} & \href{https://github.com/alectionary/vibe-math/blob/3ae21c0719cedbd928f4f7882f9b9e266567605e/formalization/distance-preservation/truncated12/catalog-chunks/F11_030.lean}{source}\\
truncated12/catalog-chunks/F11\_031.lean & \href{https://github.com/alectionary/vibe-math/blob/master/formalization/distance-preservation/truncated12/catalog-chunks/F11_031.lean}{source} & \href{https://github.com/alectionary/vibe-math/blob/3ae21c0719cedbd928f4f7882f9b9e266567605e/formalization/distance-preservation/truncated12/catalog-chunks/F11_031.lean}{source}\\
truncated12/catalog-chunks/F11\_032.lean & \href{https://github.com/alectionary/vibe-math/blob/master/formalization/distance-preservation/truncated12/catalog-chunks/F11_032.lean}{source} & \href{https://github.com/alectionary/vibe-math/blob/3ae21c0719cedbd928f4f7882f9b9e266567605e/formalization/distance-preservation/truncated12/catalog-chunks/F11_032.lean}{source}\\
truncated12/catalog-chunks/F11\_033.lean & \href{https://github.com/alectionary/vibe-math/blob/master/formalization/distance-preservation/truncated12/catalog-chunks/F11_033.lean}{source} & \href{https://github.com/alectionary/vibe-math/blob/3ae21c0719cedbd928f4f7882f9b9e266567605e/formalization/distance-preservation/truncated12/catalog-chunks/F11_033.lean}{source}\\
truncated12/catalog-chunks/F11\_034.lean & \href{https://github.com/alectionary/vibe-math/blob/master/formalization/distance-preservation/truncated12/catalog-chunks/F11_034.lean}{source} & \href{https://github.com/alectionary/vibe-math/blob/3ae21c0719cedbd928f4f7882f9b9e266567605e/formalization/distance-preservation/truncated12/catalog-chunks/F11_034.lean}{source}\\
truncated12/catalog-chunks/F11\_035.lean & \href{https://github.com/alectionary/vibe-math/blob/master/formalization/distance-preservation/truncated12/catalog-chunks/F11_035.lean}{source} & \href{https://github.com/alectionary/vibe-math/blob/3ae21c0719cedbd928f4f7882f9b9e266567605e/formalization/distance-preservation/truncated12/catalog-chunks/F11_035.lean}{source}\\
truncated12/catalog-chunks/F11\_036.lean & \href{https://github.com/alectionary/vibe-math/blob/master/formalization/distance-preservation/truncated12/catalog-chunks/F11_036.lean}{source} & \href{https://github.com/alectionary/vibe-math/blob/3ae21c0719cedbd928f4f7882f9b9e266567605e/formalization/distance-preservation/truncated12/catalog-chunks/F11_036.lean}{source}\\
truncated12/catalog-chunks/F11\_037.lean & \href{https://github.com/alectionary/vibe-math/blob/master/formalization/distance-preservation/truncated12/catalog-chunks/F11_037.lean}{source} & \href{https://github.com/alectionary/vibe-math/blob/3ae21c0719cedbd928f4f7882f9b9e266567605e/formalization/distance-preservation/truncated12/catalog-chunks/F11_037.lean}{source}\\
truncated12/catalog-chunks/F11\_038.lean & \href{https://github.com/alectionary/vibe-math/blob/master/formalization/distance-preservation/truncated12/catalog-chunks/F11_038.lean}{source} & \href{https://github.com/alectionary/vibe-math/blob/3ae21c0719cedbd928f4f7882f9b9e266567605e/formalization/distance-preservation/truncated12/catalog-chunks/F11_038.lean}{source}\\
truncated12/catalog-chunks/F11\_039.lean & \href{https://github.com/alectionary/vibe-math/blob/master/formalization/distance-preservation/truncated12/catalog-chunks/F11_039.lean}{source} & \href{https://github.com/alectionary/vibe-math/blob/3ae21c0719cedbd928f4f7882f9b9e266567605e/formalization/distance-preservation/truncated12/catalog-chunks/F11_039.lean}{source}\\
truncated12/catalog-chunks/F11\_040.lean & \href{https://github.com/alectionary/vibe-math/blob/master/formalization/distance-preservation/truncated12/catalog-chunks/F11_040.lean}{source} & \href{https://github.com/alectionary/vibe-math/blob/3ae21c0719cedbd928f4f7882f9b9e266567605e/formalization/distance-preservation/truncated12/catalog-chunks/F11_040.lean}{source}\\
truncated12/catalog-chunks/F11\_041.lean & \href{https://github.com/alectionary/vibe-math/blob/master/formalization/distance-preservation/truncated12/catalog-chunks/F11_041.lean}{source} & \href{https://github.com/alectionary/vibe-math/blob/3ae21c0719cedbd928f4f7882f9b9e266567605e/formalization/distance-preservation/truncated12/catalog-chunks/F11_041.lean}{source}\\
truncated12/catalog-chunks/F11\_042.lean & \href{https://github.com/alectionary/vibe-math/blob/master/formalization/distance-preservation/truncated12/catalog-chunks/F11_042.lean}{source} & \href{https://github.com/alectionary/vibe-math/blob/3ae21c0719cedbd928f4f7882f9b9e266567605e/formalization/distance-preservation/truncated12/catalog-chunks/F11_042.lean}{source}\\
truncated12/catalog-chunks/F11\_043.lean & \href{https://github.com/alectionary/vibe-math/blob/master/formalization/distance-preservation/truncated12/catalog-chunks/F11_043.lean}{source} & \href{https://github.com/alectionary/vibe-math/blob/3ae21c0719cedbd928f4f7882f9b9e266567605e/formalization/distance-preservation/truncated12/catalog-chunks/F11_043.lean}{source}\\
truncated12/catalog-chunks/F11\_044.lean & \href{https://github.com/alectionary/vibe-math/blob/master/formalization/distance-preservation/truncated12/catalog-chunks/F11_044.lean}{source} & \href{https://github.com/alectionary/vibe-math/blob/3ae21c0719cedbd928f4f7882f9b9e266567605e/formalization/distance-preservation/truncated12/catalog-chunks/F11_044.lean}{source}\\
truncated12/catalog-chunks/F11\_045.lean & \href{https://github.com/alectionary/vibe-math/blob/master/formalization/distance-preservation/truncated12/catalog-chunks/F11_045.lean}{source} & \href{https://github.com/alectionary/vibe-math/blob/3ae21c0719cedbd928f4f7882f9b9e266567605e/formalization/distance-preservation/truncated12/catalog-chunks/F11_045.lean}{source}\\
truncated12/catalog-chunks/F11\_046.lean & \href{https://github.com/alectionary/vibe-math/blob/master/formalization/distance-preservation/truncated12/catalog-chunks/F11_046.lean}{source} & \href{https://github.com/alectionary/vibe-math/blob/3ae21c0719cedbd928f4f7882f9b9e266567605e/formalization/distance-preservation/truncated12/catalog-chunks/F11_046.lean}{source}\\
truncated12/catalog-chunks/F11\_047.lean & \href{https://github.com/alectionary/vibe-math/blob/master/formalization/distance-preservation/truncated12/catalog-chunks/F11_047.lean}{source} & \href{https://github.com/alectionary/vibe-math/blob/3ae21c0719cedbd928f4f7882f9b9e266567605e/formalization/distance-preservation/truncated12/catalog-chunks/F11_047.lean}{source}\\
truncated12/catalog-chunks/F11\_048.lean & \href{https://github.com/alectionary/vibe-math/blob/master/formalization/distance-preservation/truncated12/catalog-chunks/F11_048.lean}{source} & \href{https://github.com/alectionary/vibe-math/blob/3ae21c0719cedbd928f4f7882f9b9e266567605e/formalization/distance-preservation/truncated12/catalog-chunks/F11_048.lean}{source}\\
truncated12/catalog-chunks/F11\_049.lean & \href{https://github.com/alectionary/vibe-math/blob/master/formalization/distance-preservation/truncated12/catalog-chunks/F11_049.lean}{source} & \href{https://github.com/alectionary/vibe-math/blob/3ae21c0719cedbd928f4f7882f9b9e266567605e/formalization/distance-preservation/truncated12/catalog-chunks/F11_049.lean}{source}\\
truncated12/catalog-chunks/F11\_050.lean & \href{https://github.com/alectionary/vibe-math/blob/master/formalization/distance-preservation/truncated12/catalog-chunks/F11_050.lean}{source} & \href{https://github.com/alectionary/vibe-math/blob/3ae21c0719cedbd928f4f7882f9b9e266567605e/formalization/distance-preservation/truncated12/catalog-chunks/F11_050.lean}{source}\\
truncated12/catalog-chunks/F11\_051.lean & \href{https://github.com/alectionary/vibe-math/blob/master/formalization/distance-preservation/truncated12/catalog-chunks/F11_051.lean}{source} & \href{https://github.com/alectionary/vibe-math/blob/3ae21c0719cedbd928f4f7882f9b9e266567605e/formalization/distance-preservation/truncated12/catalog-chunks/F11_051.lean}{source}\\
truncated12/catalog-chunks/F11\_052.lean & \href{https://github.com/alectionary/vibe-math/blob/master/formalization/distance-preservation/truncated12/catalog-chunks/F11_052.lean}{source} & \href{https://github.com/alectionary/vibe-math/blob/3ae21c0719cedbd928f4f7882f9b9e266567605e/formalization/distance-preservation/truncated12/catalog-chunks/F11_052.lean}{source}\\
truncated12/catalog-chunks/F11\_053.lean & \href{https://github.com/alectionary/vibe-math/blob/master/formalization/distance-preservation/truncated12/catalog-chunks/F11_053.lean}{source} & \href{https://github.com/alectionary/vibe-math/blob/3ae21c0719cedbd928f4f7882f9b9e266567605e/formalization/distance-preservation/truncated12/catalog-chunks/F11_053.lean}{source}\\
truncated12/catalog-chunks/F11\_054.lean & \href{https://github.com/alectionary/vibe-math/blob/master/formalization/distance-preservation/truncated12/catalog-chunks/F11_054.lean}{source} & \href{https://github.com/alectionary/vibe-math/blob/3ae21c0719cedbd928f4f7882f9b9e266567605e/formalization/distance-preservation/truncated12/catalog-chunks/F11_054.lean}{source}\\
truncated12/catalog-chunks/F11\_055.lean & \href{https://github.com/alectionary/vibe-math/blob/master/formalization/distance-preservation/truncated12/catalog-chunks/F11_055.lean}{source} & \href{https://github.com/alectionary/vibe-math/blob/3ae21c0719cedbd928f4f7882f9b9e266567605e/formalization/distance-preservation/truncated12/catalog-chunks/F11_055.lean}{source}\\
truncated12/catalog-chunks/F11\_056.lean & \href{https://github.com/alectionary/vibe-math/blob/master/formalization/distance-preservation/truncated12/catalog-chunks/F11_056.lean}{source} & \href{https://github.com/alectionary/vibe-math/blob/3ae21c0719cedbd928f4f7882f9b9e266567605e/formalization/distance-preservation/truncated12/catalog-chunks/F11_056.lean}{source}\\
truncated12/catalog-chunks/F11\_057.lean & \href{https://github.com/alectionary/vibe-math/blob/master/formalization/distance-preservation/truncated12/catalog-chunks/F11_057.lean}{source} & \href{https://github.com/alectionary/vibe-math/blob/3ae21c0719cedbd928f4f7882f9b9e266567605e/formalization/distance-preservation/truncated12/catalog-chunks/F11_057.lean}{source}\\
truncated12/catalog-chunks/F11\_058.lean & \href{https://github.com/alectionary/vibe-math/blob/master/formalization/distance-preservation/truncated12/catalog-chunks/F11_058.lean}{source} & \href{https://github.com/alectionary/vibe-math/blob/3ae21c0719cedbd928f4f7882f9b9e266567605e/formalization/distance-preservation/truncated12/catalog-chunks/F11_058.lean}{source}\\
truncated12/Catalog11.lean & \href{https://github.com/alectionary/vibe-math/blob/master/formalization/distance-preservation/truncated12/Catalog11.lean}{source} & \href{https://github.com/alectionary/vibe-math/blob/3ae21c0719cedbd928f4f7882f9b9e266567605e/formalization/distance-preservation/truncated12/Catalog11.lean}{source}\\
truncated12/Coverage11.lean & \href{https://github.com/alectionary/vibe-math/blob/master/formalization/distance-preservation/truncated12/Coverage11.lean}{source} & \href{https://github.com/alectionary/vibe-math/blob/3ae21c0719cedbd928f4f7882f9b9e266567605e/formalization/distance-preservation/truncated12/Coverage11.lean}{source}\\
truncated12/CatalogExclusion11.lean & \href{https://github.com/alectionary/vibe-math/blob/master/formalization/distance-preservation/truncated12/CatalogExclusion11.lean}{source} & \href{https://github.com/alectionary/vibe-math/blob/3ae21c0719cedbd928f4f7882f9b9e266567605e/formalization/distance-preservation/truncated12/CatalogExclusion11.lean}{source}\\
truncated12/SourceEleven.lean & \href{https://github.com/alectionary/vibe-math/blob/master/formalization/distance-preservation/truncated12/SourceEleven.lean}{source} & \href{https://github.com/alectionary/vibe-math/blob/3ae21c0719cedbd928f4f7882f9b9e266567605e/formalization/distance-preservation/truncated12/SourceEleven.lean}{source}\\
truncated12/PermuteSource.lean & \href{https://github.com/alectionary/vibe-math/blob/master/formalization/distance-preservation/truncated12/PermuteSource.lean}{source} & \href{https://github.com/alectionary/vibe-math/blob/3ae21c0719cedbd928f4f7882f9b9e266567605e/formalization/distance-preservation/truncated12/PermuteSource.lean}{source}\\
truncated12/ClassicalTransport.lean & \href{https://github.com/alectionary/vibe-math/blob/master/formalization/distance-preservation/truncated12/ClassicalTransport.lean}{source} & \href{https://github.com/alectionary/vibe-math/blob/3ae21c0719cedbd928f4f7882f9b9e266567605e/formalization/distance-preservation/truncated12/ClassicalTransport.lean}{source}\\
truncated12/EntryTestTransport.lean & \href{https://github.com/alectionary/vibe-math/blob/master/formalization/distance-preservation/truncated12/EntryTestTransport.lean}{source} & \href{https://github.com/alectionary/vibe-math/blob/3ae21c0719cedbd928f4f7882f9b9e266567605e/formalization/distance-preservation/truncated12/EntryTestTransport.lean}{source}\\
truncated12/GaugeWeight.lean & \href{https://github.com/alectionary/vibe-math/blob/master/formalization/distance-preservation/truncated12/GaugeWeight.lean}{source} & \href{https://github.com/alectionary/vibe-math/blob/3ae21c0719cedbd928f4f7882f9b9e266567605e/formalization/distance-preservation/truncated12/GaugeWeight.lean}{source}\\
truncated12/GroupTransport.lean & \href{https://github.com/alectionary/vibe-math/blob/master/formalization/distance-preservation/truncated12/GroupTransport.lean}{source} & \href{https://github.com/alectionary/vibe-math/blob/3ae21c0719cedbd928f4f7882f9b9e266567605e/formalization/distance-preservation/truncated12/GroupTransport.lean}{source}\\
truncated12/ProfileTransport.lean & \href{https://github.com/alectionary/vibe-math/blob/master/formalization/distance-preservation/truncated12/ProfileTransport.lean}{source} & \href{https://github.com/alectionary/vibe-math/blob/3ae21c0719cedbd928f4f7882f9b9e266567605e/formalization/distance-preservation/truncated12/ProfileTransport.lean}{source}\\
truncated12/CountCanonical.lean & \href{https://github.com/alectionary/vibe-math/blob/master/formalization/distance-preservation/truncated12/CountCanonical.lean}{source} & \href{https://github.com/alectionary/vibe-math/blob/3ae21c0719cedbd928f4f7882f9b9e266567605e/formalization/distance-preservation/truncated12/CountCanonical.lean}{source}\\
truncated12/TwelveConditional.lean & \href{https://github.com/alectionary/vibe-math/blob/master/formalization/distance-preservation/truncated12/TwelveConditional.lean}{source} & \href{https://github.com/alectionary/vibe-math/blob/3ae21c0719cedbd928f4f7882f9b9e266567605e/formalization/distance-preservation/truncated12/TwelveConditional.lean}{source}\\
truncated12/ThroughEleven.lean & \href{https://github.com/alectionary/vibe-math/blob/master/formalization/distance-preservation/truncated12/ThroughEleven.lean}{source} & \href{https://github.com/alectionary/vibe-math/blob/3ae21c0719cedbd928f4f7882f9b9e266567605e/formalization/distance-preservation/truncated12/ThroughEleven.lean}{source}\\
truncated12/DistanceMinima.lean & \href{https://github.com/alectionary/vibe-math/blob/master/formalization/distance-preservation/truncated12/DistanceMinima.lean}{source} & \href{https://github.com/alectionary/vibe-math/blob/3ae21c0719cedbd928f4f7882f9b9e266567605e/formalization/distance-preservation/truncated12/DistanceMinima.lean}{source}\\
truncated12/DistanceTwelve.lean & \href{https://github.com/alectionary/vibe-math/blob/master/formalization/distance-preservation/truncated12/DistanceTwelve.lean}{source} & \href{https://github.com/alectionary/vibe-math/blob/3ae21c0719cedbd928f4f7882f9b9e266567605e/formalization/distance-preservation/truncated12/DistanceTwelve.lean}{source}\\
truncated12/CriterionScope.lean & \href{https://github.com/alectionary/vibe-math/blob/master/formalization/distance-preservation/truncated12/CriterionScope.lean}{source} & \href{https://github.com/alectionary/vibe-math/blob/3ae21c0719cedbd928f4f7882f9b9e266567605e/formalization/distance-preservation/truncated12/CriterionScope.lean}{source}\\
\bottomrule\end{longtable}\normalsize

\end{document}
